% Options for packages loaded elsewhere
\PassOptionsToPackage{unicode}{hyperref}
\PassOptionsToPackage{hyphens}{url}
%
\documentclass[
]{article}
\usepackage{amsmath,amssymb}
\usepackage{iftex}
\ifPDFTeX
  \usepackage[T1]{fontenc}
  \usepackage[utf8]{inputenc}
  \usepackage{textcomp} % provide euro and other symbols
\else % if luatex or xetex
  \usepackage{unicode-math} % this also loads fontspec
  \defaultfontfeatures{Scale=MatchLowercase}
  \defaultfontfeatures[\rmfamily]{Ligatures=TeX,Scale=1}
\fi
\usepackage{lmodern}
\ifPDFTeX\else
  % xetex/luatex font selection
\fi
% Use upquote if available, for straight quotes in verbatim environments
\IfFileExists{upquote.sty}{\usepackage{upquote}}{}
\IfFileExists{microtype.sty}{% use microtype if available
  \usepackage[]{microtype}
  \UseMicrotypeSet[protrusion]{basicmath} % disable protrusion for tt fonts
}{}
\makeatletter
\@ifundefined{KOMAClassName}{% if non-KOMA class
  \IfFileExists{parskip.sty}{%
    \usepackage{parskip}
  }{% else
    \setlength{\parindent}{0pt}
    \setlength{\parskip}{6pt plus 2pt minus 1pt}}
}{% if KOMA class
  \KOMAoptions{parskip=half}}
\makeatother
\usepackage{xcolor}
\usepackage{longtable,booktabs,array}
\usepackage{calc} % for calculating minipage widths
% Correct order of tables after \paragraph or \subparagraph
\usepackage{etoolbox}
\makeatletter
\patchcmd\longtable{\par}{\if@noskipsec\mbox{}\fi\par}{}{}
\makeatother
% Allow footnotes in longtable head/foot
\IfFileExists{footnotehyper.sty}{\usepackage{footnotehyper}}{\usepackage{footnote}}
\makesavenoteenv{longtable}
\usepackage{graphicx}
\makeatletter
\def\maxwidth{\ifdim\Gin@nat@width>\linewidth\linewidth\else\Gin@nat@width\fi}
\def\maxheight{\ifdim\Gin@nat@height>\textheight\textheight\else\Gin@nat@height\fi}
\makeatother
% Scale images if necessary, so that they will not overflow the page
% margins by default, and it is still possible to overwrite the defaults
% using explicit options in \includegraphics[width, height, ...]{}
\setkeys{Gin}{width=\maxwidth,height=\maxheight,keepaspectratio}
% Set default figure placement to htbp
\makeatletter
\def\fps@figure{htbp}
\makeatother
\setlength{\emergencystretch}{3em} % prevent overfull lines
\providecommand{\tightlist}{%
  \setlength{\itemsep}{0pt}\setlength{\parskip}{0pt}}
\setcounter{secnumdepth}{-\maxdimen} % remove section numbering
% definitions for citeproc citations
\NewDocumentCommand\citeproctext{}{}
\NewDocumentCommand\citeproc{mm}{%
  \begingroup\def\citeproctext{#2}\cite{#1}\endgroup}
\makeatletter
 % allow citations to break across lines
 \let\@cite@ofmt\@firstofone
 % avoid brackets around text for \cite:
 \def\@biblabel#1{}
 \def\@cite#1#2{{#1\if@tempswa , #2\fi}}
\makeatother
\newlength{\cslhangindent}
\setlength{\cslhangindent}{1.5em}
\newlength{\csllabelwidth}
\setlength{\csllabelwidth}{3em}
\newenvironment{CSLReferences}[2] % #1 hanging-indent, #2 entry-spacing
 {\begin{list}{}{%
  \setlength{\itemindent}{0pt}
  \setlength{\leftmargin}{0pt}
  \setlength{\parsep}{0pt}
  % turn on hanging indent if param 1 is 1
  \ifodd #1
   \setlength{\leftmargin}{\cslhangindent}
   \setlength{\itemindent}{-1\cslhangindent}
  \fi
  % set entry spacing
  \setlength{\itemsep}{#2\baselineskip}}}
 {\end{list}}
\usepackage{calc}
\newcommand{\CSLBlock}[1]{\hfill\break\parbox[t]{\linewidth}{\strut\ignorespaces#1\strut}}
\newcommand{\CSLLeftMargin}[1]{\parbox[t]{\csllabelwidth}{\strut#1\strut}}
\newcommand{\CSLRightInline}[1]{\parbox[t]{\linewidth - \csllabelwidth}{\strut#1\strut}}
\newcommand{\CSLIndent}[1]{\hspace{\cslhangindent}#1}
\ifLuaTeX
\usepackage[bidi=basic]{babel}
\else
\usepackage[bidi=default]{babel}
\fi
\babelprovide[main,import]{british}
% get rid of language-specific shorthands (see #6817):
\let\LanguageShortHands\languageshorthands
\def\languageshorthands#1{}
% Consolidated manuscript style additions.
\usepackage[a4paper,margin=21mm]{geometry}
\usepackage{microtype}
\usepackage{booktabs,longtable,tabularx,array}
\usepackage{mathtools,bm}
\usepackage{caption}
\usepackage{fancyhdr}
\usepackage{xcolor}
\usepackage{titlesec}
\usepackage{hyperref}
\definecolor{chrononav}{HTML}{17365D}
\definecolor{chronoslate}{HTML}{394B59}
\definecolor{chronolight}{HTML}{EAF1F8}
\hypersetup{colorlinks=true,linkcolor=chrononav,citecolor=chrononav,urlcolor=chrononav}
\pagestyle{fancy}
\fancyhf{}
\fancyhead[L]{\small\color{chronoslate}Null-Relational Chronometry}
\fancyhead[R]{\small\color{chronoslate}Consolidated v0.1-v1.4}
\fancyfoot[C]{\thepage}
\renewcommand{\headrulewidth}{0.3pt}
\setlength{\parindent}{0pt}
\setlength{\parskip}{0.45em}
\setlength{\emergencystretch}{3em}
\captionsetup{font=small,labelfont=bf}
\titleformat{\section}{\Large\bfseries\color{chrononav}}{\thesection}{0.7em}{}
\titleformat{\subsection}{\large\bfseries\color{chronoslate}}{\thesubsection}{0.7em}{}
\titleformat{\subsubsection}{\normalsize\bfseries\color{chronoslate}}{\thesubsubsection}{0.7em}{}
\ifLuaTeX
  \usepackage{selnolig}  % disable illegal ligatures
\fi
\usepackage{bookmark}
\IfFileExists{xurl.sty}{\usepackage{xurl}}{} % add URL line breaks if available
\urlstyle{same}
\hypersetup{
  pdftitle={Null-Relational Chronometry},
  pdfauthor={Angus Muffatti},
  pdflang={en-GB},
  hidelinks,
  pdfcreator={LaTeX via pandoc}}

\title{Null-Relational Chronometry}
\usepackage{etoolbox}
\makeatletter
\providecommand{\subtitle}[1]{% add subtitle to \maketitle
  \apptocmd{\@title}{\par {\large #1 \par}}{}{}
}
\makeatother
\subtitle{From causal structure to universal spectral time: crossed-null
kinematics, scale locking, QCD-transmitted chronometric shear, protected
cosmology, and nonequilibrium transport}
\author{Angus Muffatti}
\date{Consolidated manuscript incorporating research notes v0.1-v1.4}

\begin{document}
\maketitle
\begin{abstract}
This manuscript consolidates the complete research programme developed
from the question of whether a photon, a null frame, or relations among
massless fields can supply the structure normally attributed to time.
The result is not a theory in which a photon has a rest-frame
perspective, nor a claim that causal order is created by matter. The
defensible thesis is narrower: causal or conformal structure may be
primitive, while the metric calibration of duration can be reconstructed
from physical spectra when all viable clocks share one local scalar
factor. The programme begins with exact crossed-null kinematics, records
the failure of an exact-null constrained action, replaces it with a
healthy soft-null effective theory, proves a universal-clock
factorisation theorem, and constructs a conventional scale-generating
matter model. It then derives an all-orders QCD threshold recursion and
the leading \(2/27\) transmission of a controlled scale-lock defect into
observable clock disagreement. A cyclic \(Z_6\) construction protects an
ultralight ratio mode, nonlinear environmental calculations show that
Earth and the Sun do not screen it, and a cosmological attractor selects
a nonzero branch. Exact action symmetry is reconciled with asymmetric
reheating by placing the asymmetry in the state rather than in a
permanent coupling. Closed-time-path, higher-loop matching, nonlinear
cascade, renormalisation-group, and Arnold-Moore-Yaffe transport audits
progressively remove several failure modes. The present frontier is the
complete electroweak/Yukawa LPM problem for \(H\leftrightarrow qD\) and,
beyond it, a reduced gauge-covariant Schwinger-Keldysh treatment. All
claims are classified as established, derived within the displayed
model, conditional, failed, or open; no absolute priority or
quantum-gravity completion is claimed.
\end{abstract}

{
\setcounter{tocdepth}{2}
\tableofcontents
}
\newpage

\section{Executive synthesis}\label{executive-synthesis}

\subsection{The question that started the
programme}\label{the-question-that-started-the-programme}

A photon follows a null curve and accumulates zero proper time between
emission and absorption. That fact does \textbf{not} endow it with an
inertial rest frame. A Lorentz boost to the speed of light is singular,
and coordinates adapted to a null generator do not become the
coordinates of a clock-bearing observer. Their stationary curves are
null, their intrinsic metric is degenerate, and no ordinary
three-dimensional rest space exists. The phrase ``from the photon's
perspective'' therefore fails at the level of observer structure, not
merely at the level of a difficult measurement.

The failure is productive. It isolates what a single null direction
lacks:

\begin{enumerate}
\def\labelenumi{\arabic{enumi}.}
\tightlist
\item
  a complementary null direction or nonparallel energy flow from which a
  timelike direction can be formed;
\item
  a relative normalisation that distinguishes common scale from
  rapidity;
\item
  a dimensionful spectral standard that can turn progression into
  measured duration;
\item
  an observable relative phase rather than an unobservable overall
  phase.
\end{enumerate}

The research programme asks whether those missing structures can be
generated relationally. The mature question is not ``does a photon
experience the Universe instantaneously?'' It is:

\begin{quote}
Given causal/conformal structure and interacting physical fields, under
what conditions does one operational proper-time metric emerge for every
material clock, and what observable measures the failure of that
universality?
\end{quote}

The answer developed here is:

\[
\boxed{
\text{causal order is input; universal duration calibration can be a derived field relation.}
}
\]

This is compatible with the Ehlers-Pirani-Schild programme,
causal-reconstruction theorems, and the fact that null propagation
determines conformal rather than metric geometry
\citeproc{ref-EPS1972}{{[}1{]}}, \citeproc{ref-HKM1976}{{[}2{]}},
\citeproc{ref-Malament1977}{{[}3{]}}, \citeproc{ref-Braun2025}{{[}4{]}}.
The new work lies, if anywhere, in the restricted conjunction of a
universal spectral theorem, its QCD-transmitted obstruction, a protected
cosmological realisation, and a controlled nonequilibrium preparation
mechanism. It does \textbf{not} lie in the elementary observation that
two null vectors can be combined into a timelike one.

\begin{figure}
\centering
\includegraphics[width=0.95\textwidth,height=\textheight]{figures/chronometric_emergence_chain.png}
\caption{The consolidated emergence chain.}\label{fig:emergence}
\end{figure}

\subsection{The integrated claim}\label{the-integrated-claim}

Let \((\mathcal M,[g])\) be a Lorentzian conformal spacetime. Let every
physically viable clock transition \(A\) be represented by a positive
local angular frequency \(\omega_A(x)\) of Weyl weight \(-1\), so its
accumulated phase is

\[
\Phi_A[\gamma]=\int_\gamma \omega_A\,ds_g.
\]

The phase is invariant under

\[
g_{\mu\nu}\rightarrow \Omega^2g_{\mu\nu},
\qquad
\omega_A\rightarrow \Omega^{-1}\omega_A.
\]

One conformal representative calibrates \textbf{all} clocks if and only
if

\[
\boxed{
\omega_A(x)=c_A\chi(x)
\quad\text{for every clock }A,
}
\]

where \(c_A\) is a constant and \(\chi\) is one shared Weyl-weight
\(-1\) scalar. Equivalently,

\[
\boxed{
\mathcal S_{AB}
\equiv d\ln\frac{\omega_A}{\omega_B}=0
\quad\text{for every pair }A,B.
}
\]

When factorisation holds, the operational metric is

\[
\boxed{
\widehat g_{\mu\nu}
=\left(\frac{\chi}{\chi_*}\right)^2 g_{\mu\nu},
}
\]

unique up to one global unit convention. The one-forms
\(\mathcal S_{AB}\) are called \textbf{chronometric shear}. They are
Weyl invariant and measure an obstruction that cannot be removed by
changing conformal frame.

A concrete matter model can lock

\[
\frac{h}{\chi},
\qquad
\frac{\Lambda_{\mathrm{QCD}}}{\chi},
\qquad
\frac{M_{\mathrm{P}}}{\chi}
\]

to constants. Exact locking yields rank-zero clock response and no
composition-dependent fifth force from the common radial mode. A
controlled threshold defect produces

\[
\boxed{
 d\ln\frac{\Lambda_3}{\chi}
 =\frac{2}{27}\,\varepsilon_\Psi\,d\theta+\text{controlled corrections},
}
\]

so clock drift and equivalence-principle violation become two views of
the same failure of universal spectral scaling. The coefficient \(2/27\)
is ordinary QCD threshold physics; its use as a transmission coefficient
for chronometric shear is the project-specific synthesis
\citeproc{ref-Chetyrkin1997}{{[}5{]}},
\citeproc{ref-DamourDonoghue2010}{{[}6{]}}.

\subsection{What survived, what failed, and what remains
open}\label{what-survived-what-failed-and-what-remains-open}

\begin{longtable}[]{@{}
  >{\raggedright\arraybackslash}p{(\columnwidth - 4\tabcolsep) * \real{0.3000}}
  >{\raggedleft\arraybackslash}p{(\columnwidth - 4\tabcolsep) * \real{0.4000}}
  >{\raggedright\arraybackslash}p{(\columnwidth - 4\tabcolsep) * \real{0.3000}}@{}}
\toprule\noalign{}
\begin{minipage}[b]{\linewidth}\raggedright
Layer
\end{minipage} & \begin{minipage}[b]{\linewidth}\raggedleft
Current status
\end{minipage} & \begin{minipage}[b]{\linewidth}\raggedright
Meaning
\end{minipage} \\
\midrule\noalign{}
\endhead
\bottomrule\noalign{}
\endlastfoot
Photon rest-frame language & \textbf{FAIL} & A null-adapted coordinate
system exists, but not a timelike observer frame. \\
Crossed-null clock/ruler algebra & \textbf{PASS} & Exact kinematics;
useful organising language, not a novelty claim. \\
Exact-null constrained field theory & \textbf{FAIL} & Its transverse
quadratic principal symbol is rank deficient in \(3+1\) dimensions. \\
Soft-null two-phase EFT & \textbf{PASS} & Positive Hamiltonian, real
subluminal modes for \(0<\epsilon<1\). \\
Universal spectral factorisation theorem & \textbf{PASS} & Exact
necessary-and-sufficient criterion for one clock metric. \\
Conventional scale-generating matter model & \textbf{PASS AS EFT} &
Two-derivative matter is healthy; gravity remains an EFT. \\
Turok 36-field microscopic origin & \textbf{OPEN} & Physical composite
positivity and unitarity remain disputed. \\
QCD lock and threshold recursion & \textbf{PASS AS WILSONIAN EFT} &
Exact functional recursion; coefficients calculable order by order. \\
Protected long-range ratio mode & \textbf{PASS CONDITIONALLY} & Requires
exact \(Z_6\), shift-spurion counting, and UV quality. \\
Earth/Sun screening & \textbf{PASS} & Finite bodies do not convert or
thin-shell-screen the benchmark field. \\
Cosmological branch selection & \textbf{PASS CONDITIONALLY} & Strong
low-\(f_a\) attractor benchmark; domain and inflation conditions
remain. \\
State-selected reheating & \textbf{PASS AS EFT} & Action symmetric,
state asymmetric; hard sector spurions avoided. \\
Original scalar preheating portal & \textbf{FAIL} & Tachyonic/resonant
production makes the kinematic argument unsafe. \\
Fermionic cascade repair & \textbf{PASS AT RESOLVED ORDER} & Pauli
limited and selector gated; full nonlinear gauge evolution remains. \\
RG-improved matching & \textbf{PASS AT RESOLVED ORDER} & Fixed-order
scale excursion cancels after operator running. \\
Direct QCD AMY transport & \textbf{PASS} & QCD redistribution is far
faster than reheaton decay. \\
Complete electroweak/Yukawa LPM & \textbf{OPEN} & Present decisive
kinetic calculation. \\
Full non-Abelian \(3+1\)D 2PI/KB & \textbf{OPEN} & Separate HPC
programme. \\
Emergent Einstein gravity & \textbf{OPEN} & Not derived. \\
Absolute novelty & \textbf{UNESTABLISHED} & Only a restricted candidate
conjunction is defensible. \\
\end{longtable}

\begin{figure}
\centering
\includegraphics[width=1\textwidth,height=\textheight]{figures/research_version_timeline.png}
\caption{The research did not proceed as a straight triumphal march.
Each version either killed a weak formulation, repaired it, or localised
the next uncertainty.}\label{fig:timeline}
\end{figure}

\subsection{Results taxonomy}\label{results-taxonomy}

Throughout this manuscript, statements are labelled by one of five
categories.

\begin{itemize}
\tightlist
\item
  \textbf{Established:} standard mathematics or physics already present
  in the literature.
\item
  \textbf{Derived:} an analytic result obtained inside the displayed
  model, with assumptions stated.
\item
  \textbf{Numerically verified:} a result checked by the supplied code
  within its declared discretisation or benchmark.
\item
  \textbf{Conditional:} dependent on a UV completion, approximation, or
  parameter regime not yet independently established.
\item
  \textbf{Open or failed:} either not calculated or explicitly ruled out
  by the audit.
\end{itemize}

This taxonomy is not cosmetic. The project repeatedly improved by
abandoning formulations that were attractive but wrong. The exact-null
action, the direct composite gravitational metric, the Planck-scale
cosmological benchmark, the naive reheaton branching ratio, and the
bosonic preheating portal are all retained in the manuscript as failures
because removing them would falsify the research history.

\subsection{Bibliographic audit and two corrected source-role
mismatches}\label{bibliographic-audit-and-two-corrected-source-role-mismatches}

The consolidated source ledger preserves the literature cited throughout
the working notes, but its metadata has been checked against the
corresponding records where possible. This exposed two late-stage
identifier-role mismatches that must not be silently laundered into the
final argument:

\begin{enumerate}
\def\labelenumi{\arabic{enumi}.}
\tightlist
\item
  arXiv:2605.22822 is \emph{Bottom-up open EFT for non-Abelian gauge
  theory with dynamical color environment}, not a paper deriving the
  scalar-Yukawa \(H\leftrightarrow qD\) LPM normalisation. It is
  relevant to the proposed gauge-covariant open-EFT direction, but it
  does not support the direct scalar conversion rate.
\item
  arXiv:2211.15454 is \emph{AMY Lorentz invariant parton cascade---the
  thermal equilibrium case}. It supports AMY cascade implementation and
  equilibrium validation, but not the narrower description ``direct
  numerical solution of the AMY transverse integral equations.''
\end{enumerate}

Accordingly, neither identifier is used here as evidentiary support for
the v1.4 \(H\leftrightarrow qD\) rate. That rate is retained as a
project-derived numerical result whose supplied consolidated
verification script checks the reported benchmark algebra, while the
\textbf{complete electroweak/Yukawa LPM derivation remains open}. The
source ledger records both the corrected metadata and the original
project role so that the provenance remains inspectable.

\section{Part I - Null frames, photons, and relational
time}\label{part-i---null-frames-photons-and-relational-time}

\subsection{1. Why a photon has no ordinary
perspective}\label{why-a-photon-has-no-ordinary-perspective}

For light in flat spacetime,

\[
c^2\Delta t^2-\Delta x^2=0,
\]

and therefore the proper-time interval along the ray is

\[
\Delta\tau=0.
\]

This statement is invariant. It does not mean that the emission and
absorption events are one spacetime event, and it does not define an
experience of instantaneous traversal. A Lorentz transformation to a
frame with \(v=c\) does not exist. If an observer chases the ray at
\(v<c\), then for a coordinate interval \(T\) in the laboratory frame,

\[
\Delta t'
=\gamma T(1-\beta)
=T\sqrt{\frac{1-\beta}{1+\beta}},
\]

which tends to zero as \(\beta\rightarrow1\). The limiting values exist;
the limiting frame does not. At \(\beta=1\) the boost is noninvertible.

A useful way to expose the degeneration is to introduce photon-adapted
coordinates

\[
T=t,
\qquad
\xi=x-ct.
\]

The ray sits at \(\xi=0\), but the metric becomes

\[
 ds^2=d\xi^2+2c\,dT\,d\xi+dy^2+dz^2,
\]

with \(g_{TT}=0\). Curves of constant \((\xi,y,z)\) are null. They
cannot carry functioning proper-time clocks. The coordinates are valid;
the putative observer frame is not.

The deeper obstruction is geometric. A massive observer has timelike
four-velocity \(U^\mu\) with \(U^2=-1\), and the subspace orthogonal to
\(U\) is a spacelike rest space. A null tangent \(k^\mu\) satisfies
\(k^2=0\) and belongs to its own orthogonal complement. The quotient

\[
Q_k=k^\perp/\langle k\rangle
\]

is a two-dimensional screen space, not an observer's three-dimensional
rest space. Null tetrads and Carrollian geometry are the appropriate
structures \citeproc{ref-Gourgoulhon2006}{{[}7{]}}.

This gives the first sober ontology:

\begin{quote}
A single photon is better represented as a null connection between
interactions than as a small enduring object carrying an internal clock.
\end{quote}

\subsection{2. Compactifying observer
space}\label{compactifying-observer-space}

Massive unit four-velocities form a hyperboloid. Write a boost in
direction \(\hat{\mathbf n}\) as

\[
U_\eta=(\cosh\eta,\sinh\eta\,\hat{\mathbf n}).
\]

The vector diverges as \(\eta\rightarrow\infty\), but the projectively
rescaled limit exists:

\[
\lim_{\eta\to\infty}2e^{-\eta}U_\eta
=(1,\hat{\mathbf n})
\equiv k,
\qquad k^2=0.
\]

The ideal photon frame is therefore a projective null ray \([k]\), not a
unit four-velocity. Massive observers occupy the interior of hyperbolic
velocity space; null directions form its ideal boundary. This is a
geometric completion, not an algebraic invention analogous to adjoining
\(i\) to the reals.

An ideal null frame can be represented as

\[
\mathcal F_\gamma=(\gamma,[k],Q_k,\mathcal A_\gamma),
\]

where \(\gamma\) is a null geodesic, \(Q_k\) is its screen space, and
\(\mathcal A_\gamma\) is its affine structure with

\[
\lambda\sim a\lambda+b.
\]

It retains causal order, transverse geometry, polarisation, and affine
ratios. It does not contain proper time, longitudinal metric distance,
an absolute energy scale, or ordinary simultaneity.

\subsection{3. Two null directions and the timelike
interior}\label{two-null-directions-and-the-timelike-interior}

Let \(k^a\) and \(\ell^a\) be future-directed, nonparallel null
covectors with

\[
k^2=\ell^2=0,
\qquad
k\cdot\ell=-2C<0.
\]

Define

\[
u^a=\frac{k^a+\ell^a}{2\sqrt C},
\qquad
e^a=\frac{k^a-\ell^a}{2\sqrt C}.
\]

Then

\[
u^2=-1,
\qquad
e^2=+1,
\qquad
u\cdot e=0.
\]

This is standard null-basis algebra. Its interpretive value is that
reciprocal rescaling

\[
k\rightarrow e^\eta k,
\qquad
\ell\rightarrow e^{-\eta}\ell
\]

acts as

\[
\begin{aligned}
 u&\rightarrow\cosh\eta\,u+\sinh\eta\,e,\\
 e&\rightarrow\sinh\eta\,u+\cosh\eta\,e,
\end{aligned}
\]

so rapidity is the relative logarithmic weighting of two null
directions.

Independent rescalings \(k\rightarrow ak\), \(\ell\rightarrow b\ell\)
separate into

\[
\alpha=\frac12\ln(ab),
\qquad
\eta=\frac12\ln\frac ab.
\]

The ratio controls boost; the product controls common scale. Two
unscaled null rays do not select a unique observer. Physical momenta do,
because their energies fix the weighting. For two nonparallel null
momenta \(p\) and \(q\),

\[
P^a=p^a+q^a,
\qquad
M^2=-P^2=-2p\cdot q>0,
\]

and

\[
u^a=\frac{P^a}{M}
\]

is their centre-of-momentum time direction. An ensemble of nonparallel
null momenta generally lies inside the convex future cone and admits a
timelike Landau energy frame. A single plane wave does not.

The useful slogan is therefore precise but limited:

\[
\boxed{
\text{timelike structure occupies the convex interior generated by null relations.}
}
\]

It does not yet fix a unit of time.

\subsection{4. Why a single light phase cannot be its own
clock}\label{why-a-single-light-phase-cannot-be-its-own-clock}

For a plane wave with phase \(\phi(x)=k_\mu x^\mu\) and a ray

\[
x^\mu(\lambda)=x_0^\mu+\lambda k^\mu,
\]

one has

\[
\frac{d\phi}{d\lambda}=k_\mu k^\mu=0.
\]

The phase is constant along its own generator. Oscillation appears to a
timelike detector, not as an internal pendulum attached to the ray. An
autonomous clock requires comparison: two nonparallel modes, a cavity, a
bound state, emission and absorption events, or some other interaction
that supplies a second phase and a scale.

This observation motivates the crossed-phase construction.

\section{Part II - Crossed-null
chronometry}\label{part-ii---crossed-null-chronometry}

\subsection{5. Exact crossed-phase
theorem}\label{exact-crossed-phase-theorem}

Let \(\phi_+\) and \(\phi_-\) be circle-valued phases with null
gradients

\[
k_a=\nabla_a\phi_+,
\qquad
\ell_a=\nabla_a\phi_-,
\]

\[
g^{ab}k_ak_b=g^{ab}\ell_a\ell_b=0.
\]

Assume they are nonparallel and define

\[
C=-\frac12g^{ab}k_a\ell_b>0.
\]

Set

\[
T=\frac{\phi_++\phi_-}{2},
\qquad
R=\frac{\phi_+-\phi_-}{2},
\]

and introduce the conformally invariant tensor

\[
h_{ab}=C g_{ab}.
\]

Because

\[
\nabla T=\frac{k+\ell}{2},
\qquad
\nabla R=\frac{k-\ell}{2},
\]

the inverse Gram matrix under \(g\) is

\[
\begin{pmatrix}
\nabla T\cdot\nabla T & \nabla T\cdot\nabla R\\
\nabla R\cdot\nabla T & \nabla R\cdot\nabla R
\end{pmatrix}
=
\begin{pmatrix}
-C&0\\
0&C
\end{pmatrix}.
\]

Since \(h^{ab}=C^{-1}g^{ab}\),

\[
\boxed{
\begin{aligned}
h^{ab}\nabla_aT\nabla_bT&=-1,\\
h^{ab}\nabla_aR\nabla_bR&=+1,\\
h^{ab}\nabla_aT\nabla_bR&=0.
\end{aligned}}
\]

Thus two crossed null phases define a unit clock direction and a unit
ruler direction in their two-plane. Under
\(g_{ab}\rightarrow\Omega^2g_{ab}\),

\[
C\rightarrow\Omega^{-2}C,
\]

so \(h_{ab}=Cg_{ab}\) is invariant.

In flat spacetime, choose

\[
\phi_+=\omega_+(t-x),
\qquad
\phi_-=\omega_-(t+x).
\]

Then

\[
C=\omega_+\omega_-,
\qquad
\eta=\frac12\ln\frac{\omega_+}{\omega_-}.
\]

Therefore

\[
\boxed{
\omega_+\omega_- = \text{clock or mass scale}^{2},
\qquad
\frac{\omega_+}{\omega_-}=\text{rapidity data}.
}
\]

The wave sum factorises as

\[
e^{i\phi_+}+e^{i\phi_-}=2e^{iT}\cos R.
\]

The mean phase is a temporal carrier and the difference phase supplies
standing-wave nodes. This is the algebraic skeleton of a clock-and-ruler
pair.

\subsubsection{Novelty correction}\label{novelty-correction}

The exact map is not a new conformal completion. Writing
\(\phi_\pm=T\pm R\) and imposing exact nullity gives

\[
C=-(\nabla T)^2,
\]

and hence

\[
h_{\mu\nu}=-(g^{\alpha\beta}\partial_\alpha T\partial_\beta T)g_{\mu\nu},
\]

which is the singular conformal map of mimetic gravity, with \(R\)
supplying an orthogonal equal-norm ruler field
\citeproc{ref-Mimetic2013}{{[}8{]}}. The theorem remains exact and
useful; the broad novelty claim does not.

\subsection{6. Homogeneous locking potential and its
limitation}\label{homogeneous-locking-potential-and-its-limitation}

The first proposed mechanism introduced a mass-squared order parameter
\(\Sigma\) and

\[
V(C,\Sigma)
=\frac\kappa2(C-\Sigma)^2
+\frac\beta4\Sigma^2\left(\ln\frac{\Sigma}{\Lambda^2}-\frac12\right)
+\frac\beta8\Lambda^4.
\]

The stationary point

\[
C=\Sigma=\Lambda^2
\]

has Hessian

\[
H=
\begin{pmatrix}
\kappa&-\kappa\\
-\kappa&\kappa+\beta/2
\end{pmatrix},
\]

so

\[
H_{11}=\kappa>0,
\qquad
\det H=\frac{\kappa\beta}{2}>0.
\]

For \(\kappa,\beta>0\), this is a strict homogeneous minimum. It
demonstrates a possible algebraic lock of a crossed phase invariant to a
transmuted scale. It does \textbf{not} establish kinetic rank,
hyperbolicity, absence of ghosts, or a healthy \(3+1\) dimensional field
theory.

That distinction proved decisive.

\subsection{7. Failure of the exact-null
action}\label{failure-of-the-exact-null-action}

The first complete constrained action was

\[
S_0=\int d^4x\sqrt{-g}\left[
\frac{M_{\mathrm{P}}^2}{2}R
+A_+Y_++A_-Y_-
-\frac\kappa2(C-\sigma^2)^2
-\frac12(\nabla\sigma)^2
-V(\sigma)
\right],
\]

where

\[
Y_\pm=(\nabla\phi_\pm)^2,
\qquad
C=-\frac12\nabla\phi_+\cdot\nabla\phi_-.
\]

Around

\[
\bar\phi_+=m(t-x),
\qquad
\bar\phi_-=m(t+x),
\]

the linearised constraints are

\[
(\partial_t+\partial_x)\pi_+=0,
\qquad
(\partial_t-\partial_x)\pi_-=0.
\]

Define

\[
s=\frac{\pi_++\pi_-}{\sqrt2},
\qquad
r=\frac{\pi_+-\pi_-}{\sqrt2}.
\]

The lock fluctuation is

\[
\delta C=\frac{m}{\sqrt2}(\dot s-\partial_x r).
\]

Neither the linear constraints nor the quadratic lock contains
\(\partial_y\) or \(\partial_z\). Transverse gradients first enter
nonlinearly through terms such as \((\nabla\pi_\pm)^2\). The quadratic
principal symbol is therefore rank deficient for generic transverse
perturbations. Nonlinear terms control directions with no quadratic
restoring operator.

\[
\boxed{
\text{Verdict: the exact-null constrained action is not a controlled ordinary }3+1\text{D EFT.}
}
\]

The positive homogeneous Hessian did not detect this because it tested
only algebraic order-parameter fluctuations.

\subsection{8. Healthy soft-null repair}\label{healthy-soft-null-repair}

The minimal repair makes the phases ordinary propagating compact fields
and treats crossed nullity as a background state:

\[
\boxed{
\mathcal L_{\mathrm{EFT}}
=-\frac F2(Y_++Y_-)
+\frac\kappa2(C-m^2)^2.
}
\]

Define

\[
a=\frac{\kappa m^2}{2},
\qquad
\epsilon=\frac aF.
\]

The quadratic Lagrangian in \(s,r\) is

\[
\begin{aligned}
\mathcal L^{(2)}={}&
\frac12(F+a)\dot s^2
+\frac12F\dot r^2
-a\dot s\,\partial_xr\\
&-\frac F2[(\nabla s)^2+(\nabla r)^2]
+\frac a2(\partial_xr)^2.
\end{aligned}
\]

The Hamiltonian is positive when

\[
\boxed{F>0,\qquad0<a<F,}
\]

or \(0<\epsilon<1\). The principal polynomial factorises:

\[
F(\omega^2-k^2)
\left[(F+a)\omega^2-Fk^2+ak_x^2\right]=0.
\]

The modes are

\[
\omega_1^2=k^2,
\]

\[
\omega_2^2=\frac{Fk^2-ak_x^2}{F+a}.
\]

Writing \(k_x=k\cos\theta\),

\[
\boxed{
c_2^2(\theta)=\frac{1-\epsilon\cos^2\theta}{1+\epsilon}.
}
\]

For \(0<\epsilon<1\), the second mode is real, positive, and subluminal
for every direction. The model belongs to the established
two-superfluid/entrainment class rather than defining a new species of
matter \citeproc{ref-Superfluid2015}{{[}9{]}}.

\subsubsection{Local mediator
completion}\label{local-mediator-completion}

Introduce a healthy massive scalar \(H\):

\[
\mathcal L_{\mathrm{med}}
=-\frac F2(Y_++Y_-)
-\frac12(\nabla H)^2
-\frac12M^2H^2
+g_HH(C-m^2).
\]

Integrating it out below \(M\) gives

\[
\Delta\mathcal L
=\frac{g_H^2}{2}(C-m^2)\frac1{M^2-\Box}(C-m^2),
\]

and therefore

\[
\boxed{\kappa=\frac{g_H^2}{M^2}>0.}
\]

The sign is compatible with exchange of a positive-residue massive state
and with standard positivity reasoning
\citeproc{ref-Adams2006}{{[}10{]}}.

\subsubsection{Model-specific tree-level sum
rules}\label{model-specific-tree-level-sum-rules}

Parallel and perpendicular speeds satisfy

\[
c_\parallel^2=\frac{1-\epsilon}{1+\epsilon},
\qquad
c_\perp^2=\frac1{1+\epsilon},
\]

hence

\[
\boxed{c_\parallel^2=2c_\perp^2-1.}
\]

For the minimal one-mediator completion, the \(k^4\) dispersion
coefficients also obey

\[
\boxed{\alpha_\parallel=4\alpha_\perp.}
\]

These are not universal laws of two-fluid physics. They are consistency
relations of the minimal model and will be perturbed by additional
operators, resonances, and loops.

\subsection{9. Gravity coupling and
circularity}\label{gravity-coupling-and-circularity}

The direct substitution

\[
g^{\mathrm{physical}}_{\mu\nu}=Cg_{\mu\nu}
\]

is singular and mimetic. It does not evade the noninvertibility or
extra-mode problem. A regular alternative introduces a Weyl compensator
\(\chi\):

\[
g_{\mu\nu}\rightarrow\Omega^2g_{\mu\nu},
\qquad
\chi\rightarrow\Omega^{-1}\chi,
\]

\[
\widehat g_{\mu\nu}
=\frac{\chi^2}{6M_{\mathrm{P}}^2}g_{\mu\nu}.
\]

A representative action is

\[
\begin{aligned}
S_W=\int\sqrt{-g}\Big[&
\frac{\chi^2}{12}R
+\frac12(\nabla\chi)^2
-\frac{f^2\chi^2}{2}(Y_++Y_-)\\
&+\frac\kappa2(C-\zeta\chi^2)^2
\Big].
\end{aligned}
\]

On the locking surface, the phase state calibrates the same invariant
metric used by gravity. This is a regular EFT representation of
chronometric calibration. It does not derive Einstein dynamics from null
phases. The conformal manifold and gravitational action remain inputs.

\section{Part III - Universal spectral
chronometry}\label{part-iii---universal-spectral-chronometry}

\subsection{10. Clock spectral fields}\label{clock-spectral-fields}

A local clock is represented not by a coordinate label but by a positive
transition frequency \(\omega_A(x)\). The invariant quantity is its
phase accumulation

\[
d\Phi_A=\omega_A ds_g.
\]

Each individual clock defines a representative

\[
g^{(A)}_{\mu\nu}
=\left(\frac{\omega_A}{\omega_{A*}}\right)^2g_{\mu\nu},
\]

for which its intrinsic rate is constant. The nontrivial question is
whether every clock defines the same representative.

\subsection{11. Universal clock factorisation
theorem}\label{universal-clock-factorisation-theorem}

\textbf{Theorem.} Let \(U\subset\mathcal M\) be connected and let
\(\{\omega_A\}\) be positive clock spectral fields. The following are
equivalent:

\begin{enumerate}
\def\labelenumi{\arabic{enumi}.}
\tightlist
\item
  There exists \(\widehat g\in[g]\) and positive constants
  \(\bar\omega_A\) such that
  \(\Phi_A[\gamma]=\bar\omega_A\int_\gamma ds_{\widehat g}\) for every
  timelike curve \(\gamma\) in \(U\).
\item
  Every pairwise ratio is constant: \(d\ln(\omega_A/\omega_B)=0\).
\item
  There is one positive Weyl-weight \(-1\) scalar \(\chi\) and constants
  \(c_A\) such that \(\omega_A=c_A\chi\).
\end{enumerate}

\textbf{Proof.} Write \(\widehat g=e^{2\sigma}g\). Equality of phase
integrands gives \(\omega_A=\bar\omega_Ae^\sigma\), proving constant
ratios. If all ratios are constant, choose one reference clock and set
\(\chi\) proportional to its spectrum; the other spectra are constant
multiples. Conversely, if \(\omega_A=c_A\chi\), then
\(\widehat g=(\chi/\chi_*)^2g\) gives

\[
\Phi_A=c_A\chi_*\int ds_{\widehat g}.
\]

Uniqueness is up to one global multiplicative unit convention.
\(\square\)

The theorem is elementary once stated. Its value is that it turns the
phrase ``all clocks measure the same proper time'' into an exact
integrability condition.

\subsection{12. Chronometric shear and clock-space
rank}\label{chronometric-shear-and-clock-space-rank}

Define

\[
\boxed{
\mathcal S_{AB}=d\ln\frac{\omega_A}{\omega_B}.
}
\]

This is invariant under Weyl rescaling. A useful clock-space
decomposition begins with

\[
\ell_A=\ln\omega_A,
\]

where conformal gauge shifts all components in the common direction
\((1,\ldots,1)\). Physical clock disagreement lives in the quotient of
log-spectrum space by that common direction.

With positive weights \(w_A\) satisfying \(\sum_Aw_A=1\), define

\[
B^A_\mu=-\partial_\mu\ln\omega_A,
\qquad
W_\mu=\sum_Aw_AB^A_\mu,
\]

\[
\Sigma^A_\mu=B^A_\mu-W_\mu.
\]

Then \(\Sigma^A\) and \(\mathcal S_{AB}\) are gauge invariant, and

\[
\mathfrak C_{\mu\nu}=\sum_Aw_A\Sigma^A_\mu\Sigma^A_\nu
\]

is positive semidefinite. It vanishes exactly when one chronometric
metric exists.

Suppose spectra depend on \(r\) nonuniversal dimensionless fields
\(q^I\):

\[
\delta\ln\omega_A
=\delta\ln\chi+K_{AI}\delta q^I.
\]

Clock ratios obey

\[
\delta\ln\frac{\omega_A}{\omega_B}
=(K_{AI}-K_{BI})\delta q^I,
\]

so

\[
\boxed{
\operatorname{rank}\left[\delta\ln(\omega_A/\omega_B)\right]\le r.
}
\]

One nonuniversal scalar predicts a rank-one network across arbitrarily
many clocks. Exact universal chronometry is rank zero.

\subsection{13. The role of
electromagnetism}\label{the-role-of-electromagnetism}

Source-free Maxwell theory in four dimensions is conformally invariant.
It transports phase, null structure, polarisation, and causal
information, but it does not select an absolute spectral scale. A free
wave's frequency is state data.

Matter interactions create bound-state gaps. A hydrogenic scale is
schematically

\[
E_{\mathrm{Ry}}\sim m_e\alpha_{\mathrm{em}}^2,
\qquad
m_e=\frac{y_eh}{\sqrt2}.
\]

If \(E_A=c_A\chi\), then

\[
d\Phi_A=\frac{E_A}{\hbar}ds_g
=c_A\chi_*ds_{\widehat g}.
\]

The division of labour is therefore:

\[
\boxed{
\begin{aligned}
\chi&:\text{ universal normalisation},\\
\text{Higgs and QCD}&:\text{ transmission into material spectra},\\
\text{electromagnetism}&:\text{ phase organisation, transport, and readout}.
\end{aligned}}
\]

The correct statement is not ``the Higgs and photons create time.'' It
is that a deeper scale field calibrates duration, material sectors relay
the scale, and electromagnetic phase makes the count operational.

\subsection{14. Why the Standard Model Higgs is insufficient by
itself}\label{why-the-standard-model-higgs-is-insufficient-by-itself}

A general clock frequency has the structure

\[
\omega_A
=h\,F_A\left(
\alpha_{\mathrm{em}},g_s,y_f,
\frac{\Lambda_{\mathrm{QCD}}}h,
\frac{M_{\mathrm{P}}}h,\ldots
\right).
\]

The ordinary Higgs sets elementary electroweak masses, but it does not
impose

\[
d\ln\frac{\Lambda_{\mathrm{QCD}}}h=0,
\qquad
d\ln\frac{M_{\mathrm{P}}}h=0.
\]

Nuclear clocks depend on confinement, quark masses, nuclear magnetic
moments, and many-body structure. The Higgs is a relay, not
automatically a universal source.

Higgs-only induced gravity would require

\[
M_{\mathrm{P}}^2=\xi_hv^2,
\]

so

\[
\xi_h\sim10^{32},
\]

far beyond the minimal phenomenologically plausible range
\citeproc{ref-Atkins2012}{{[}11{]}}. Using the sole radial Higgs mode as
a local Weyl compensator also removes the physical scalar remaining
after the electroweak Goldstones are eaten. A second radial direction is
structurally required.

\subsection{15. Higgs-dilaton and deeper-scale
completion}\label{higgs-dilaton-and-deeper-scale-completion}

Introduce a positive scalar \(\chi\) and impose

\[
\boxed{
h=\alpha\chi,
\qquad
\Lambda_{\mathrm{QCD}}=c_Q\chi,
\qquad
M_{\mathrm{P}}=c_P\chi.
}
\]

If all dimensionless couplings are fixed, every physical gap is

\[
\omega_A=c_A\chi.
\]

Write

\[
h=\rho\sin\theta,
\qquad
\chi=\rho\cos\theta.
\]

The radial coordinate \(\rho\) is the common scale; the angular
coordinate \(\theta\) measures Higgs/dilaton alignment. A pure radial
fluctuation rescales every clock together and is invisible to local
ratios. An angular fluctuation produces

\[
\delta\ln\frac{\omega_A}{\omega_B}
=(K_A-K_B)\delta\theta.
\]

A healthy globally scale-invariant representative has

\[
F(h,\chi)=\xi_hh^2+\xi_\chi\chi^2,
\]

\[
V(h,\chi)=\frac\lambda4(h^2-\alpha^2\chi^2)^2+\beta\chi^4.
\]

In the Einstein frame, the scalar field-space metric

\[
G_{ij}=\frac{M_{\mathrm{P}}^2}{F}\delta_{ij}
+\frac{3M_{\mathrm{P}}^2}{2F^2}F_{,i}F_{,j}
\]

is positive for positive Jordan kinetic signs and \(F>0\). In the exact
global scale limit, one radial dilaton direction can remain flat; local
Weyl symmetry can gauge it; quantum breaking can lift it. These are
different physical theories and should not be conflated
\citeproc{ref-Shaposhnikov2009}{{[}12{]}},
\citeproc{ref-Foot2007}{{[}13{]}}.

\subsection{16. A conventional hidden-sector
generator}\label{a-conventional-hidden-sector-generator}

A safe infrared proof of principle uses

\[
G=G_{\mathrm{SM}}\times SU(2)_X,
\]

with the Standard Model Higgs doublet \(H\) and a hidden scalar doublet
\(\Phi\):

\[
h=\sqrt{2H^\dagger H},
\qquad
\chi=\sqrt{2\Phi^\dagger\Phi}.
\]

The Jordan-frame action includes

\[
\frac12(\xi_\chi\chi^2+\xi_Hh^2)R,
\]

ordinary two-derivative kinetic terms, the hidden gauge field, and

\[
V
=\frac{\lambda_\chi}{2}(\Phi^\dagger\Phi)^2
-\lambda_p(H^\dagger H)(\Phi^\dagger\Phi)
+\frac{\lambda_H}{2}(H^\dagger H)^2.
\]

Hidden gauge loops generate a Coleman-Weinberg scale
\(f=\langle\chi\rangle\). The portal valley gives

\[
h^2=\frac{\lambda_p}{\lambda_H}\chi^2
\equiv\zeta_h^2\chi^2,
\]

and the gravitational coefficient becomes

\[
F=\xi_{\mathrm{eff}}\chi^2,
\qquad
M_{\mathrm{P}}^2=\xi_{\mathrm{eff}}f^2.
\]

This construction is not proposed as the final UV theory. Its role is to
demonstrate that a conventional, unitary, two-derivative matter sector
can generate and transmit one common scale
\citeproc{ref-HiddenSU22013}{{[}14{]}},
\citeproc{ref-CurvedSM2018}{{[}15{]}}.

\subsubsection{Exact common mode and fifth-force
cancellation}\label{exact-common-mode-and-fifth-force-cancellation}

In the Einstein frame,

\[
m_A^E=\frac{M_{\mathrm{P}}^{(0)}m_A(\chi)}{\sqrt{F(\chi)}}.
\]

If \(m_A=c_A\chi\) and \(F=\xi_{\mathrm{eff}}\chi^2\), then

\[
m_A^E=\frac{c_AM_{\mathrm{P}}^{(0)}}{\sqrt{\xi_{\mathrm{eff}}}},
\]

independent of the radial field. Therefore

\[
\alpha_A\equiv\partial_\varphi\ln m_A^E=0.
\]

Exact universal scaling produces no composition-dependent fifth force.
Imperfect locking gives

\[
\alpha_A-\alpha_B
=\partial_\varphi\ln\frac{m_A}{m_B}
=\partial_\varphi\ln\frac{\omega_A}{\omega_B}.
\]

Differential fifth-force charge and differential clock drift are the
same derivative obstruction
\citeproc{ref-ScaleFifthForce2016}{{[}16{]}},
\citeproc{ref-ScaleFifthForce2021}{{[}17{]}}.

\subsection{17. Turok's 36 fields as an optional microscopic
origin}\label{turoks-36-fields-as-an-optional-microscopic-origin}

Boyle and Turok proposed 36 dimension-zero fourth-order scalar fields in
a programme involving anomaly cancellation, vacuum energy, and an
emergent Higgs \citeproc{ref-BoyleTurok2022}{{[}18{]}}. A possible
chronometric connection is that the 36-field sector might generate one
collective infrared singlet \(\chi\), for example through a mass gap or
an invariant composite related to

\[
\mathcal O_2
=\frac1{36}G_{IJ}g^{\mu\nu}
\nabla_\mu\varphi^I\nabla_\nu\varphi^J.
\]

The mapping would be

\[
36\text{ microscopic dimension-zero fields}
\longrightarrow
1\text{ dimensionful infrared order parameter}.
\]

However, an elementary propagator with leading behaviour \(1/p^4\)
cannot possess a nontrivial conventional positive Källén-Lehmann
spectral measure. If

\[
G(Q^2)=\int_0^\infty ds\,\frac{\rho(s)}{Q^2+s},
\qquad\rho(s)\ge0,
\]

then the leading term is

\[
G(Q^2)=\frac1{Q^2}\int\rho(s)ds+O(Q^{-4}).
\]

To begin at \(1/Q^4\) requires \(\int\rho=0\), which with positivity
forces \(\rho=0\). This does not exclude an indefinite-metric formalism,
a gauge redundancy, or a healthy composite pole. It means physical
positivity must be demonstrated for invariant observables, not inferred
from the elementary fourth-order field.

Bateman and Turok propose a Krein-space and hidden-ghost-parity
resolution, while Cline and Hell argue that ghosts, unitarity violation,
and a confining fifth force remain
\citeproc{ref-BatemanTurok2026}{{[}19{]}},
\citeproc{ref-ClineHell2026}{{[}20{]}}. The chronometric programme
therefore treats the 36-field theory as an optional UV candidate subject
to four mandatory tests:

\begin{enumerate}
\def\labelenumi{\arabic{enumi}.}
\tightlist
\item
  Lorentzian unitarity on a physical state space;
\item
  positive spectral density or an equally strong reconstruction theorem
  for an invariant composite;
\item
  a stable condensate with \(Z_\chi>0\) and real dispersion;
\item
  acceptable universal matter coupling without a macroscopic confining
  force.
\end{enumerate}

The safe low-energy construction survives whether or not that candidate
passes.

\section{Part IV - QCD locking and the first nonzero
prediction}\label{part-iv---qcd-locking-and-the-first-nonzero-prediction}

\subsection{18. Scale-invariant quantum
definition}\label{scale-invariant-quantum-definition}

To prevent QCD from acquiring an independent subtraction standard, work
in \(d=4-2\varepsilon\) and replace the external regularisation scale by

\[
\mu_\chi=z\,\chi^{1/(1-\varepsilon)},
\]

where \(z\) is dimensionless. The bare strong coupling is

\[
g_{s,0}=\mu_\chi^\varepsilon Z_g g_s.
\]

Renormalisation generates an infinite tower of scale-invariant
higher-dimensional operators. The result is an all-orders Wilsonian EFT,
not a finite-parameter UV completion. In a homogeneous background, every
physical gap takes the form

\[
M_A=f\,\mathcal F_A(g_i,z,c_k,\ldots),
\]

and

\[
\Lambda_{n_f}=c_{n_f}\chi.
\]

Slowly varying backgrounds receive derivative corrections such as

\[
\Lambda_{n_f}(x)=c_{n_f}\chi(x)
\left[1+c_\Box\frac{\Box\chi}{\chi^3}
+c_\partial\frac{(\nabla\chi)^2}{\chi^4}+\cdots\right].
\]

The physical statement is not merely ``evaluate \(g_s\) at
\(\mu=\chi\).'' It is that the quantum theory contains no independent
dimensionful standard besides \(\chi\).

\subsection{19. Exact threshold propagation
theorem}\label{exact-threshold-propagation-theorem}

Let a heavy coloured particle \(j\) with renormalisation-group-invariant
mass \(\widehat m_j\) be integrated out. Dimensional analysis gives

\[
\Lambda_{n-1}
=\widehat m_j\,
\mathcal F_j\left(\frac{\Lambda_n}{\widehat m_j}\right).
\]

Define

\[
A_j=\frac{\partial\ln\mathcal F_j}
{\partial\ln(\Lambda_n/\widehat m_j)}.
\]

Differentiation yields

\[
d\ln\Lambda_{n-1}
=A_jd\ln\Lambda_n
+(1-A_j)d\ln\widehat m_j.
\]

Define the lock and threshold defects

\[
\delta_n=d\ln\frac{\Lambda_n}{\chi},
\qquad
\epsilon_j=d\ln\frac{\widehat m_j}{\chi}.
\]

Then

\[
\boxed{
\delta_{n-1}=A_j\delta_n+(1-A_j)\epsilon_j.
}
\]

For many thresholds,

\[
\delta_3
=\left(\prod_{n=4}^{N}A_n\right)\delta_N
+\sum_{n=4}^{N}(1-A_n)
\left(\prod_{k=4}^{n-1}A_k\right)\epsilon_n.
\]

If the high-energy theory and every threshold are locked, the low-energy
theory remains locked to all orders. Scheme and loop dependence live in
the physical response coefficients \(A_j\), not in the structure of the
recursion \citeproc{ref-Chetyrkin1997}{{[}5{]}}.

\subsection{20. A controlled coloured
defect}\label{a-controlled-coloured-defect}

Introduce one real scalar \(S\) and a vectorlike Dirac fermion

\[
\Psi\sim(\mathbf3,\mathbf1,0),
\]

with

\[
\mathcal L_{\Psi S}
=\bar\Psi(i\gamma^\mu D_\mu-y_\chi\chi-y_SS)\Psi
-\frac12(\partial S)^2
-V_{\chi S},
\]

\[
V_{\chi S}
=V_\chi(\chi)
+\frac{\lambda_S}{4}(S^2-r^2\chi^2)^2
+\Delta V_{\mathrm{ct}}.
\]

At alignment, \(S=r\chi\) and

\[
M_\Psi=\chi(y_\chi+y_Sr).
\]

Parameterise the ratio fluctuation by

\[
\frac S\chi=re^\theta.
\]

Then

\[
\boxed{
 d\ln\frac{M_\Psi}{\chi}
 =\varepsilon_\Psi d\theta,
\qquad
\varepsilon_\Psi=\frac{y_Sr}{y_\chi+y_Sr}.
}
\]

This is a physical crack in universal locking: it changes a
dimensionless threshold ratio rather than merely rescaling every mass.

\subsection{\texorpdfstring{21. The \(2/27\) transmission
coefficient}{21. The 2/27 transmission coefficient}}\label{the-227-transmission-coefficient}

The threshold chain is

\[
7\xrightarrow{\Psi}6\xrightarrow{t}5\xrightarrow{b}4\xrightarrow{c}3.
\]

Assume the high-energy QCD scale and the top, bottom, and charm
thresholds remain exactly locked. At one loop,

\[
A_\Psi=\frac{19}{21},
\qquad
A_t=\frac{21}{23},
\qquad
A_b=\frac{23}{25},
\qquad
A_c=\frac{25}{27}.
\]

Therefore

\[
\begin{aligned}
\mathcal D_{\Psi\to3}^{(1)}
&=\left(1-\frac{19}{21}\right)
\frac{21}{23}\frac{23}{25}\frac{25}{27}\\
&=\boxed{\frac{2}{27}}.
\end{aligned}
\]

The first controlled defect is

\[
\boxed{
 d\ln\frac{\Lambda_3}{\chi}
 =\frac{2}{27}\varepsilon_\Psi d\theta
 +O\left(\alpha_s\varepsilon_\Psi,
 \varepsilon_\Psi^2,
 \frac{\partial^2}{\chi^2}\right).
}
\]

For a Dirac fermion in colour representation \(R\),

\[
\mathcal D_{R\to3}^{(1)}=\frac{4T(R)}{27}.
\]

Again, the coefficient is not new. The new use is to make the first
low-energy failure of universal time calibration calculable rather than
arbitrary.

\subsection{22. Clock and equivalence-principle
predictions}\label{clock-and-equivalence-principle-predictions}

A low-energy clock ratio obeys

\[
\mathcal S_{AB}
=\Delta K_\alpha d\ln\alpha
+\Delta K_\mu d\ln\mu
+\Delta K_q d\ln X_q+\cdots,
\]

with

\[
\mu=\frac{m_p}{m_e},
\qquad
X_q=\frac{\widehat m}{\Lambda_3}.
\]

In the pure-QCD benchmark,

\[
d\ln\alpha=0,
\qquad
d\ln\mu\simeq\delta_3,
\qquad
d\ln X_q=-\delta_3.
\]

Thus

\[
\boxed{
\mathcal S_{AB}
=(\Delta K_\mu-\Delta K_q)
\frac{2}{27}\varepsilon_\Psi d\theta+\cdots.
}
\]

At leading sensitivity, the illustrative ratios give

\[
\mathcal S_{\mathrm{Sr}/Cs}
=2\mathcal S_{\mathrm{Sr}/CaF}.
\]

This is a rank-one pattern, conditional on one dominant scalar and the
stated clock sensitivities \citeproc{ref-ClockData2023}{{[}21{]}},
\citeproc{ref-ClockReview2025}{{[}22{]}}.

Let \(a=f_a\theta\) and \(\varphi=a/M_{\mathrm{P}}\). The gluonic
coupling is

\[
d_g
=\frac{\partial}{\partial\varphi}
\ln\frac{\Lambda_3}{\chi}
=\frac{2}{27}\varepsilon_\Psi\frac{M_{\mathrm{P}}}{f_a}+\cdots.
\]

In the one-scalar, unscreened, QCD-dominated limit, the anomalous clock
redshift \(\beta_{AB}\) and Eötvös parameter \(\eta_{CD}\) satisfy

\[
\boxed{
\frac{\beta_{AB}}{\eta_{CD}}
=-\frac{\Delta K_\mu-\Delta K_q}
{\Delta Q'_{\widehat m}{}^{CD}}.
}
\]

For the project's leading Ti/Pt charge model, the conditional values are
approximately \(150.38\) for Sr/CaF and \(300.76\) for Sr/Cs. They are
model-consistency lines, not new experimental limits
\citeproc{ref-DamourDonoghue2010}{{[}6{]}},
\citeproc{ref-MICROSCOPE2022}{{[}23{]}}.

\subsection{23. The radiative naturalness
obstruction}\label{the-radiative-naturalness-obstruction}

The coloured fermion contributes

\[
\Delta V_\Psi
=-\frac{N_c}{16\pi^2}M_\Psi(\theta)^4
\left[\ln\frac{M_\Psi(\theta)^2}{\mu^2}-\frac32\right].
\]

The ratio mode receives

\[
\boxed{
m_a^2\sim\frac{N_c}{4\pi^2}
\varepsilon_\Psi^2\frac{M_\Psi^4}{f_a^2}.
}
\]

A long-range observable signal is therefore catastrophically unnatural
in the minimal one-threshold model. The programme requires a symmetry
that suppresses the vacuum mass more strongly than the visible linear
coupling.

\section{Part V - Cyclic protection and environmental
response}\label{part-v---cyclic-protection-and-environmental-response}

\subsection{\texorpdfstring{24. Exact \(Z_N\)
orbit}{24. Exact Z\_N orbit}}\label{exact-z_n-orbit}

Take

\[
\mathcal G
=\left[\prod_{k=0}^{N-1}SU(3)_k\right]\rtimes Z_N,
\]

with one vectorlike colour triplet \(\Psi_k\) in each sector and a
compact field \(a/f_a\). The masses form the exact orbit

\[
\boxed{
M_k(a)=M\left[1-\epsilon\cos\left(\frac a{f_a}+\frac{2\pi k}{N}\right)\right].
}
\]

All masses are positive for \(0<\epsilon<1\). The full vacuum sums all
sectors, but visible experiments probe only sector zero.

The root-of-unity identity ensures that for \(m<N\),

\[
\sum_{k=0}^{N-1}\cos^m\left(x+\frac{2\pi k}{N}\right)
\]

is independent of \(x\). The first nonconstant Coleman-Weinberg harmonic
is therefore the \(N\)th:

\[
V_{\mathrm{CW}}(a)
=\frac{N_cM^4\epsilon^N}{8\pi^2}F_N
\cos\frac{Na}{f_a}[1+O(\epsilon)],
\]

where for \(N>4\),

\[
F_N=\frac{24\,2^{1-N}}
{(N-1)(N-2)(N-3)(N-4)}.
\]

The one-sector visible threshold derivative remains \(O(\epsilon)\),
while the complete-orbit vacuum mass is \(O(\epsilon^N)\). This
established discrete-Goldstone mechanism is adapted here to a real
vectorlike QCD threshold \citeproc{ref-DasHook2020}{{[}24{]}},
\citeproc{ref-Brzeminski2020}{{[}25{]}},
\citeproc{ref-DiLuzio2021}{{[}26{]}}.

\begin{figure}
\centering
\includegraphics[width=0.88\textwidth,height=\textheight]{figures/z6_mass_orbit_and_potential.png}
\caption{The complete orbit cancels lower vacuum harmonics while
preserving a one-sector linear threshold response.}\label{fig:z6}
\end{figure}

\subsection{\texorpdfstring{25. Minimal clean choice:
\(Z_6\)}{25. Minimal clean choice: Z\_6}}\label{minimal-clean-choice-z_6}

The smallest even \(N>4\) that supplies a symmetry-related vacuum with
maximal visible slope is \(N=6\). Choose

\[
x_0=\frac{a_0}{f_a}=\frac\pi2.
\]

Then \(M_0=M\) and

\[
\frac{\partial}{\partial x}
\ln\frac{M_0}{\chi}=\epsilon.
\]

Since \(F_6=1/160\),

\[
\boxed{
m_a^2
=\frac{27}{320\pi^2}
\frac{M^4}{f_a^2}\epsilon^6+O(\epsilon^7).
}
\]

The visible QCD coupling remains

\[
\boxed{
d_g=\frac{2}{27}\frac{M_{\mathrm{P}}}{f_a}\epsilon
+O(\alpha_s\epsilon,\epsilon^2).
}
\]

Compared with the unprotected mass,

\[
\boxed{
\frac{m_{a,Z_6}^2}{m_{a,\mathrm{naive}}^2}
=\frac9{80}\epsilon^4.
}
\]

This makes an astronomical-range scalar compatible with a much larger
visible response. The quality conditions are severe: the exact exchange
symmetry must be microscopic, and local operators such as \(\Phi^6\)
must be forbidden or nonlocally suppressed. A Wilson-line or
collective-pNGB completion is the preferred route
\citeproc{ref-Hor2026}{{[}27{]}}.

\subsection{26. Nonlinear environmental
equation}\label{nonlinear-environmental-equation}

Let \(x=a/f_a\). The static spherical equation is

\[
\nabla^2x
=-\frac{m_a^2}{6}\sin6x
+\sum_A\frac{\rho_A}{f_a^2}
 p_A\epsilon\sin x
 (1-\epsilon\cos x)^{p_A-1}.
\]

The benchmark

\[
M=10\,\mathrm{TeV},
\qquad
f_a=M_{\mathrm{P}},
\qquad
\epsilon=10^{-6}
\]

has

\[
m_a=3.797\times10^{-21}\,\mathrm{eV},
\qquad
\lambda_a=347.4\,\mathrm{AU}.
\]

\subsubsection{Homogeneous spinodal versus finite-body
conversion}\label{homogeneous-spinodal-versus-finite-body-conversion}

Define the local density parameter

\[
t=\frac{\rho p\epsilon}{m_a^2f_a^2}.
\]

The branch connected to \(x_\infty=\pi/2\) disappears in an infinite
homogeneous medium above

\[
t_*=0.17272340294,
\]

corresponding in the benchmark to

\[
\rho_{\mathrm{spinodal}}=46.25\,\mathrm{g\,cm^{-3}}.
\]

The solar core exceeds this local threshold. A finite body must also pay
gradient energy. For a core of radius \(R_c\), define

\[
q(R_c)=\frac{M(<R_c)p\epsilon}{4\pi f_a^2R_c}.
\]

Conversion to the adjacent phase requires at least

\[
\boxed{q(R_c)>q_{\mathrm{conv}}=0.633135.}
\]

Earth, Sun, and a 5 cm tungsten source fall short by enormous margins.

\begin{longtable}[]{@{}
  >{\raggedright\arraybackslash}p{(\columnwidth - 6\tabcolsep) * \real{0.2000}}
  >{\raggedleft\arraybackslash}p{(\columnwidth - 6\tabcolsep) * \real{0.2667}}
  >{\raggedleft\arraybackslash}p{(\columnwidth - 6\tabcolsep) * \real{0.2667}}
  >{\raggedleft\arraybackslash}p{(\columnwidth - 6\tabcolsep) * \real{0.2667}}@{}}
\toprule\noalign{}
\begin{minipage}[b]{\linewidth}\raggedright
Source
\end{minipage} & \begin{minipage}[b]{\linewidth}\raggedleft
Largest \(q/q_{\mathrm{conv}}\)
\end{minipage} & \begin{minipage}[b]{\linewidth}\raggedleft
Surface shift
\end{minipage} & \begin{minipage}[b]{\linewidth}\raggedleft
Screening factor
\end{minipage} \\
\midrule\noalign{}
\endhead
\bottomrule\noalign{}
\endlastfoot
Earth & \(1.63\times10^{-16}\) & \(-1.03\times10^{-16}\) &
\(1.000000\) \\
Sun & \(1.02\times10^{-12}\) & \(-3.15\times10^{-13}\) & \(1.000000\) \\
5 cm tungsten & \(3.50\times10^{-32}\) & \(-2.22\times10^{-32}\) &
\(1.000000\) \\
\end{longtable}

\begin{figure}
\centering
\includegraphics[width=0.82\textwidth,height=\textheight]{figures/environmental_conversion_margin.png}
\caption{Density can prefer another phase locally, but finite objects
are much too weakly compact to realise it.}\label{fig:environment}
\end{figure}

The Sun therefore provides the desired source without screening it. The
field's response length is vastly larger than the source, so it reacts
to integrated scalar charge rather than tracking a local density
minimum. This is a thick-shell/no-shell regime rather than a chameleon
thin shell \citeproc{ref-Chameleon2008}{{[}28{]}},
\citeproc{ref-FiniteDensity2024}{{[}29{]}}.

\subsubsection{Finite-size bound}\label{finite-size-bound}

Using \(G=(8\pi M_{\mathrm{P}}^2)^{-1}\),

\[
q=2p\epsilon\left(\frac{M_{\mathrm{P}}}{f_a}\right)^2\Phi,
\]

where \(\Phi=GM/(Rc^2)\). If

\[
f_a\ge M_{\mathrm{P}},
\qquad
\epsilon\le1,
\qquad
p\le\frac{2}{27},
\]

then any nonsingular spherical body outside a trapped surface has

\[
q<\frac{2}{27}<q_{\mathrm{conv}}.
\]

This excludes complete matter-driven conversion for that parameter
class, not cosmological domains or more elaborate UV interactions.

\section{Part VI - Cosmological branch
selection}\label{part-vi---cosmological-branch-selection}

\subsection{27. Homogeneous evolution}\label{homogeneous-evolution}

The cosmological mode obeys

\[
\ddot x+3H\dot x+\frac{1}{f_a^2}
\frac{\partial V_{\mathrm{eff}}(x,t)}{\partial x}=0,
\]

with

\[
V_{\mathrm{eff}}=V_0+V_{\Psi,T}+V_{{\mathrm{QCD}},T}+V_b.
\]

The protected vacuum is

\[
V_0(x)=\frac{m_a^2f_a^2}{N^2}(1+\cos Nx),
\]

with minima \(x_v=(2j+1)\pi/N\).

\subsection{28. Reheating phasor}\label{reheating-phasor}

Let sector temperatures be \(T_k=\xi_kT_0\). The leading field-dependent
high-temperature threshold potential is governed by

\[
W_2=\sum_{k=0}^{N-1}\xi_k^2e^{2\pi ik/N},
\]

and selects

\[
x_T=-\arg W_2.
\]

The successful sparse population is

\[
\xi_0=1,
\qquad
\xi_5=0.25,
\qquad
\xi_{1,2,3,4}\ll1.
\]

For \(N=6\),

\[
x_T=0.0524383.
\]

The bias scales as \(\xi^2\), while dark radiation scales as \(\xi^4\).
One adjacent complete Standard Model copy at \(\xi=0.25\) gives the
project value

\[
\Delta N_{\mathrm{eff}}=0.0289,
\]

whereas populating all five hidden replicas at that temperature would be
excessive under the cited cosmological bound
\citeproc{ref-GoldsteinHill2026}{{[}30{]}}.

\subsection{29. QCD crossover as a second
lens}\label{qcd-crossover-as-a-second-lens}

Write the QCD pressure as

\[
p(T,\Lambda)=T^4P(T/\Lambda).
\]

Dimensional analysis gives

\[
\frac{\partial F}{\partial\ln\Lambda}=\rho-3p\equiv\Theta(T).
\]

Since

\[
\Lambda_k(x)=\Lambda_0
\left[1-\epsilon\cos\left(x+\frac{2\pi k}{N}\right)\right]^{2/27},
\]

the leading QCD term is

\[
\boxed{
V_{{\mathrm{QCD}},k}
=\frac{2}{27}\Theta(T_k)
\ln\left[1-\epsilon\cos\left(x+\frac{2\pi k}{N}\right)\right].
}
\]

The lattice-QCD trace anomaly produces a second focusing epoch near
\(T\simeq0.179\) GeV in the strong benchmark
\citeproc{ref-Borsanyi2016}{{[}31{]}}.

\subsection{30. Baryon locking and late
release}\label{baryon-locking-and-late-release}

Visible baryons contribute

\[
V_b(x,a)=\rho_b(a)(1-\epsilon\cos x)^{2/27}.
\]

Near the even-\(N\) thermal point, baryons supply positive curvature.
The vacuum potential eventually becomes tachyonic there and releases the
field into an adjacent minimum. In the preferred benchmark, this occurs
around

\[
z_{\mathrm{inst}}\simeq450.
\]

The chronology is

\[
\boxed{
\text{heavy-threshold focusing}
\rightarrow
\text{QCD refocusing}
\rightarrow
\text{baryon locking}
\rightarrow
\text{late vacuum release}.
}
\]

\subsection{31. Strong-attractor benchmark and
ridge}\label{strong-attractor-benchmark-and-ridge}

The original Planck-scale benchmark was cosmologically unsatisfactory:
thermal focusing was negligible and generic misalignment overclosed the
Universe. The successful strong-attractor point is

\[
\boxed{
\begin{aligned}
N&=6,\\
f_a&=2.435\times10^{10}\,\mathrm{GeV},\\
\epsilon&=2.70\times10^{-13},\\
M&=1.002\times10^6\,\mathrm{GeV},\\
T_R&=1.002\times10^8\,\mathrm{GeV},\\
m_a&=7.5\times10^{-29}\,\mathrm{eV},\\
d_g&=10^{-6}.
\end{aligned}}
\]

The focusing measures are

\[
\eta_{\Psi,\max}=612,
\qquad
\eta_{{\mathrm{QCD}},\max}=759.
\]

Twelve tested homogeneous initial phases converged to

\[
x_v=+\frac\pi6.
\]

This is a numerical existence result, not a global theorem. A moderate
point with \(\eta_{\Psi,\max}\approx30.6\) did not select a unique
branch. Merely having \(m_T/H>1\) is insufficient.

The viable \(N=6\) ridge at representative \(m_a\) and the stated cuts
is approximately

\[
1.6\times10^9\,\mathrm{GeV}
\lesssim f_a\lesssim
9.4\times10^{12}\,\mathrm{GeV},
\]

\[
2.6\times10^3\,\mathrm{GeV}
\lesssim M\lesssim
1.6\times10^7\,\mathrm{GeV}.
\]

Cosmology chooses a viable branch and sign; equivalence-principle tests
cap the magnitude of shear.

\subsection{32. Observable scale and domain
conditions}\label{observable-scale-and-domain-conditions}

For annual solar-potential modulation
\(\Delta\Phi_\odot=3.29\times10^{-10}\),

\[
\left|\Delta\ln\frac{\nu_A}{\nu_B}\right|_{\mathrm{yr}}
=2|K_{AB}|d_g^2\Delta\Phi_\odot.
\]

At \(d_g=10^{-6}\),

\[
\boxed{
\left|\Delta\ln\frac{\nu_A}{\nu_B}\right|_{\mathrm{yr}}
=6.58\times10^{-22}|K_{AB}|.
}
\]

Ordinary atomic sensitivities are far below current reach. Provisional
thorium strong-sector enhancements of order \(10^4\) would raise the
illustrative signal to \(6.58\times10^{-18}\), but the nuclear
coefficients remain model dependent
\citeproc{ref-Arakawa2026}{{[}32{]}},
\citeproc{ref-Thorium2026}{{[}33{]}}.

The six vacua require either pre-inflation symmetry breaking with no
restoration or a gauged/discrete completion. A sufficient homogeneity
condition is

\[
\frac{H_{\mathrm{inf}}}{2\pi f_a}\ll x_T.
\]

For the benchmark,

\[
H_{\mathrm{inf}}\ll8.0\times10^9\,\mathrm{GeV}.
\]

The compact field is not the dark matter on the optimised ridge; its
abundance is negligible.

\section{Part VII - Exact symmetry with an asymmetric cosmological
state}\label{part-vii---exact-symmetry-with-an-asymmetric-cosmological-state}

\subsection{33. Why hard asymmetric reheating is
dangerous}\label{why-hard-asymmetric-reheating-is-dangerous}

The cyclic cancellation of lower vacuum harmonics relies on exact
microscopic symmetry. A permanent sector-specific reheating coupling is
a hard \(Z_6\)-breaking spurion and generically feeds lower harmonics
back into the vacuum potential. The correct architecture is therefore

\[
\boxed{
\text{exact }Z_6\text{ action}
+\text{temporarily asymmetric state}.
}
\]

This distinction is strongly supported by later cyclic finite-density
models \citeproc{ref-Delaunay2026}{{[}34{]}}.

\subsection{34. Orbit-symmetric reheaton and transient
selector}\label{orbit-symmetric-reheaton-and-transient-selector}

Introduce two complete reheaton orbits \(X_k,Y_k\) with oriented cyclic
mixing. Light eigenstates have the schematic form

\[
R_k=\cos\theta\,X_k+\sin\theta\,Y_{k-1}.
\]

A selected population of \(R_0\) decays into sectors 0 and 5. A complete
selector orbit \(Q_k\) has one-hot inflationary branches. On the \(Q_0\)
branch, all reheaton copies except \(R_0\) are made kinematically
inaccessible by an exactly \(Z_6\)-invariant mass pattern.

The key chronology is

\[
\Gamma_{\phi\rightarrow R_0R_0}
\gg\Gamma_Q\gg\Gamma_R.
\]

The inflaton first produces the labelled state. The selector returns to
the symmetric origin before the reheaton decays. The asymmetry survives
in occupation numbers, not in late-time masses or couplings.

At the first homogeneous level, a branching ratio
\(\Gamma_5/\Gamma_0=1/256\) gives \(T_5/T_0=1/4\). The nonlinear cascade
later corrects the exact mixing required to achieve that final
temperature ratio.

\subsection{35. Perturbations and
isocurvature}\label{perturbations-and-isocurvature}

For fixed branching fractions,

\[
\rho_i=B_i\rho_R,
\]

and radiation curvature perturbations satisfy

\[
\zeta_i=\zeta_R+\frac14\delta\ln B_i.
\]

If the branching fractions are fixed constants,

\[
\delta B_0=\delta B_5=0,
\]

so

\[
S_{50}=3(\zeta_5-\zeta_0)=0.
\]

The extra radiation is adiabatic at linear order. This would fail if the
mixing angle were a light modulus, because its fluctuation would be
amplified by the small adjacent branching fraction. The benchmark keeps
the mixing and selector sectors heavy during inflation.

The initial chronometric phase fluctuation

\[
\delta x_{\mathrm{reh}}=\frac{H_*}{2\pi f_a}
\]

is strongly damped by threshold, QCD, and baryon focusing. The staged
calculation found a transfer magnitude of about \(10^{-6}\) by
recombination and a residual QCD-lock fluctuation near \(10^{-23}\).
This remains a reduced perturbation calculation, not a full
Einstein-Boltzmann likelihood.

\section{Part VIII - Closed-time-path radiative
audit}\label{part-viii---closed-time-path-radiative-audit}

\subsection{36. Vacuum, selector, and state
functionals}\label{vacuum-selector-and-state-functionals}

The in-in effective action separates schematically as

\[
\Gamma_{\mathrm{CTP}}
=\Gamma_{\mathrm{vac}}+\Gamma_Q+\Gamma_\rho.
\]

Exact microscopic \(Z_6\) implies

\[
\Gamma_{\mathrm{vac}}[a]
=\sum_{q\in\mathbb Z}c_{6q}e^{i6qa/f_a},
\]

so

\[
c_p^{\mathrm{vac}}=0,
\qquad p=1,\ldots,5.
\]

During a transient one-hot selector background, the action may contain

\[
\epsilon^p e^{ipa/f_a}\mathcal Q_{-p}[Q]+\mathrm{h.c.},
\]

where the selector carries compensating cyclic charge. These terms
vanish when \(Q\rightarrow0\).

An asymmetric density matrix can also produce lower harmonics through
occupation charges

\[
\mathcal N_p=\sum_kn_ke^{-2\pi ipk/6}.
\]

These are intended finite-density forces. They dilute or disappear with
the state. The central covariance relation is

\[
\Gamma_{z\rho z^{-1}}[z\phi,zQ]
=\Gamma_\rho[\phi,Q].
\]

A forbidden lower harmonic must therefore be accompanied by selector or
state charge. Under the displayed interaction graph, neither can become
a permanent state-independent Wilson coefficient merely by disappearing.

\subsection{37. Two-loop result}\label{two-loop-result}

At two loops, the selector-reheaton-bath graph and the \(a\)-dependent
threshold-QCD graph are not directly joined by a primitive vacuum
vertex. Selector information reaches the threshold through propagators
and occupation numbers prepared by reheating. Decomposing

\[
S_k=S_{k,\mathrm{vac}}+\delta S_{k,\rho},
\qquad
D_k=D_{k,\mathrm{vac}}+\delta D_{k,\rho},
\]

shows that the complete-orbit all-vacuum term cancels the forbidden
harmonics; every \(p<6\) contribution contains at least one state
insertion.

The physical equation takes the form

\[
\Box a+V_0'(a)
+\sum_kM_k'(a)\langle\bar\Psi_k\Psi_k\rangle_{\rho,\mathrm{ren}}=0.
\]

In influence-functional language, the environment supplies:

\begin{itemize}
\tightlist
\item
  a local finite-density force;
\item
  a retarded memory and dissipation kernel;
\item
  a noise kernel.
\end{itemize}

A decaying state can leave temporary memory but not automatically a
permanent local vacuum term. A conserved relic or hydrodynamic zero mode
would preserve a state-dependent source, not a new vacuum spurion.

The NLO thermal-QCD audit increased focusing by about \(6.19\%\) and
shifted the selected phase by only \(1.72\times10^{-4}\) rad. The branch
mechanism was perturbatively stable at this order
\citeproc{ref-ThermalPressure2021}{{[}35{]}}.

\subsection{38. Initial-state
renormalisation}\label{initial-state-renormalisation}

Non-vacuum initial states can introduce ultraviolet divergences. For
UV-soft or Hadamard-like data, ordinary bulk counterterms suffice. More
general short-distance structure requires counterterms on the initial
hypersurface,

\[
S_{\partial,\rho}
=\int_{t=t_0}d^3x\,\mathcal L_{\partial,\rho}.
\]

This is distinct from a permanent late-time bulk term. The preparation
protocol must still be explicitly renormalised
\citeproc{ref-CollinsHolman2005}{{[}36{]}},
\citeproc{ref-Berges2005}{{[}37{]}}.

\subsection{39. Necessary protection
package}\label{necessary-protection-package}

Exact \(Z_6\) is necessary but not sufficient. The UV theory must also
satisfy

\[
\epsilon\rightarrow0
\quad\Longrightarrow\quad
U(1)_a\text{ shift symmetry is restored}.
\]

Every nonderivative \(a\) dependence must carry explicit powers of the
same breaking spurion. The robust package is

\[
\boxed{
\text{exact }Z_6
+\text{continuous-shift spurion counting}
+\text{no direct }Q\text{-threshold portal}
+\text{UV-soft state preparation}.
}
\]

\section{Part IX - Mixed higher loops, preheating repair, and
Wilson-line
quality}\label{part-ix---mixed-higher-loops-preheating-repair-and-wilson-line-quality}

\subsection{40. First mixed selector-threshold
topology}\label{first-mixed-selector-threshold-topology}

Using only the displayed interactions

\[
Q^\dagger QR^2,
\qquad
RH^\dagger H,
\qquad
H\bar qD,
\qquad
a\bar DD,
\]

the first connected mixed graph has

\[
I=8,
\qquad
V=6,
\qquad
L=I-V+1=3.
\]

It is one-particle irreducible but two-particle reducible. In the
self-consistent 2PI description, it is generated by coupled lower-loop
kernels rather than one primitive 2PI skeleton. The transient operator
has

\[
\Delta V_{Qa}^{(3)}
=\epsilon C_3
\operatorname{Re}[e^{ia/f_a}\mathcal Q_1]
+O(\epsilon^2).
\]

A first NDA estimate made it negligible relative to the intended thermal
potential. The subsequent exact factorised matching confirmed that
conclusion.

\subsection{41. Two-time surrogate and what it did not
prove}\label{two-time-surrogate-and-what-it-did-not-prove}

The project evolved complete unequal-time propagators for an exact
quadratic non-Markovian surrogate with 20 coordinates, 10 momentum
modes, 1,801 time steps, and 301 stored two-time slices. It preserved
the symplectic structure to \(7.11\times10^{-15}\), reproduced
\(T_5/T_0=1/4\), and suppressed unselected leakage to approximately
\(1.14\times10^{-7}\).

This was a genuine two-time consistency test. It was not a nonlinear,
non-Abelian \(3+1\) dimensional plasma simulation. The distinction
remains explicit \citeproc{ref-Tranberg2024}{{[}38{]}},
\citeproc{ref-Lindner2007}{{[}39{]}},
\citeproc{ref-Bhattacharya2022}{{[}40{]}}.

\subsection{42. Failure of the direct bosonic inflaton
portal}\label{failure-of-the-direct-bosonic-inflaton-portal}

The original trilinear

\[
\mathcal L_{\phi R}
=-\frac{g_{\phi R}}2\phi R^2
\]

requires a coupling large enough that post-inflationary oscillations
make selected and nominally heavy reheaton modes tachyonic. The
illustrative resonance parameter was
\(q_{\mathrm{end}}\sim6\times10^3\). Late-time kinematic closure
therefore does not prevent early tachyonic or resonant production. This
portal fails generically \citeproc{ref-Dufaux2006}{{[}41{]}}.

\subsection{43. Selector-gated fermionic
cascade}\label{selector-gated-fermionic-cascade}

Replace it with a complete fermionic parent orbit \(N_k\):

\[
\mathcal L_N
=-y_\phi\phi\sum_k\bar N_kN_k
-\sum_km_{N_k}(Q)\bar N_kN_k
-y_R\sum_kR_k\bar N_k\nu_k+\mathrm{h.c.}
\]

On the one-hot branch, \(N_0\) is accessible while the other copies are
extremely heavy. The hidden Landau-Zener exponent in the benchmark is
about \(460.4\), giving a linear production suppression near
\(10^{-200}\). Selected fermions can be nonperturbatively produced, but
Pauli blocking caps their occupancy rather than allowing unbounded
bosonic amplification \citeproc{ref-GreeneKofman1998}{{[}42{]}}.

The repair passes at the resolved linear and cascade levels. Full
backreaction, rescattering, and gauge-plasma production remain part of
the nonlinear programme.

\subsection{44. Deconstructed Wilson-line
completion}\label{deconstructed-wilson-line-completion}

A candidate quality completion uses an 18-link deconstructed gauge
theory. The collective Wilson line is

\[
\mathcal W
=\prod_{j=0}^{17}\frac{\Sigma_j}{f/\sqrt2}
=e^{ia/f_a}.
\]

A cyclic transformation shifts replica labels and rotates
\(\mathcal W\). A messenger chain generates

\[
\mathcal L_{\Psi,\mathrm{eff}}
=-\sum_k\left[
M-\kappa(\omega^k\mathcal W+\omega^{-k}\mathcal W^\dagger)
\right]\bar\Psi_k\Psi_k.
\]

The first local winding operator requires all 18 links, so low-dimension
local operators cannot directly depend on the complete phase. This is
the type of nonlocal/collective protection needed to turn the infrared
\(Z_6\) model into a credible gauge-quality completion
\citeproc{ref-Hor2026}{{[}27{]}}.

Perturbative discrete anomaly sums pass for the displayed vectorlike
orbits, but the full global/cobordism anomaly problem is open
\citeproc{ref-Byakti2017}{{[}43{]}}.

\section{Part X - Exact matching, nonlinear cascade, and RG
completion}\label{part-x---exact-matching-nonlinear-cascade-and-rg-completion}

\subsection{45. Factorised three-loop
matching}\label{factorised-three-loop-matching}

Because the first mixed graph is two-particle reducible, its
zero-momentum coefficient factorises into a one-loop selector kernel and
a mass derivative of the known two-loop fermion-fermion-scalar effective
potential \citeproc{ref-Martin2001}{{[}44{]}}. In the scalar proxy,

\[
\mathcal I_3(\bar\mu)
=2m_R^2K_{RRH}(m_R^2,m_h^2)
D_{FFS}(M^2,m_h^2;\bar\mu).
\]

At \(\bar\mu=M\),

\[
\boxed{
\mathcal I_3(M)=6.57973508149\simeq\frac{2\pi^2}{3}.
}
\]

The first proxy coefficient gave

\[
|\Delta V_{Qa}^{(3)}|
=2.1723\times10^3\,\mathrm{GeV}^4,
\]

or only \(1.546\times10^{-12}\) of the intended thermal focusing
potential.

The fixed-order hard function had a huge apparent scale excursion. This
was explicitly identified as missing RG completion, not a physical
uncertainty.

\subsection{46. Spurion-graded operator
mixing}\label{spurion-graded-operator-mixing}

A correction to the earlier claim was necessary. Exact \(Z_6\) does not
make the complete invariant transient basis block diagonal, because
combinations \(e^{ipx}\mathcal X_{-p}\) are individually invariant. The
correct selection rule is

\[
\boxed{
\gamma_{(A,p)(B,q)}
=\sum_{r,s\ge0}
\delta^{(6)}_{p-q,r-s}
\epsilon^{r+s}\Gamma_{AB}^{(r,s)}.
}
\]

Changing harmonic charge requires explicit shift-breaking spurions. The
powers of \(\epsilon\) are fixed, although the finite tensors depend on
the messenger completion.

\subsection{47. Nonlinear momentum-lattice cascade and corrected
branching}\label{nonlinear-momentum-lattice-cascade-and-corrected-branching}

The repaired sequence is

\[
\phi\rightarrow N_0\bar N_0
\rightarrow R_0\nu_0
\rightarrow H_0,H_5
\rightarrow D_k,q_k,g_k.
\]

A radial momentum-lattice calculation included expansion, backreaction,
two-body decays, Pauli blocking, and an energy-conserving plasma
closure. It exposed a genuine error in the v0.9 branching argument.

The massless \(\nu_0\) daughter deposits visible-sector energy earlier
than the reheaton products and experiences a different redshift history.
Therefore the attractive exact choice

\[
\tan\theta=\frac1{16}
\]

does not yield the desired final temperature ratio. The simulation gave

\[
\mathcal R_\nu
=\frac{E_\nu}{E_R^{(1)}}
=0.35551328.
\]

Imposing \(E_5/E_0=1/256\) requires

\[
\boxed{
B_5^*=\frac{1+\mathcal R_\nu}{257}
=0.00527437074,
}
\]

The v1.2 note reported

\[
\tan\theta=0.0728196.
\]

Direct recomputation from the displayed value of \(B_5^*\) gives

\[
\tan\theta
=\sqrt{\frac{B_5^*}{1-B_5^*}}
=0.07281715.
\]

The difference is a small historical rounding inconsistency and is
recorded rather than silently erased. The resulting v1.2 calculation
obtained

\[
\frac{T_5}{T_0}=0.250000001.
\]

After restoring the complete Higgs-doublet and relativistic-\(D\)
contribution to \(g_*\) in v1.3, the superseding branch was

\[
B_5=0.00529888708,
\qquad
\tan\theta=0.0729871.
\]

This correction is retained as a warning against replacing a resolved
cascade with an aesthetically neat branching identity.

\subsection{48. RG-improved hard
function}\label{rg-improved-hard-function}

The fixed-order function was

\[
\begin{aligned}
D_{FFS}^{\mathrm{fixed}}(\bar\mu)
={}&2\ln\frac xz\ln\frac{x-z}{\bar\mu^2}
-\ln^2\frac{x}{\bar\mu^2}\\
&-2\operatorname{Li}_2\left(\frac zx\right)
+\frac{\pi^2}{3}.
\end{aligned}
\]

Combining it with operator running and the matching counterterm yields

\[
\boxed{
D_{FFS}^{\mathrm{hard}}
=2\ln\frac xz\ln\left(1-\frac zx\right)
-2\operatorname{Li}_2\left(\frac zx\right)
+\frac{\pi^2}{3}.
}
\]

Its explicit matching-scale derivative vanishes. The benchmark remains

\[
\mathcal I_3^{\mathrm{hard}}=6.57973508149.
\]

Ordinary coupling and portal running leaves only

\[
0.96469<\frac{C_3(\mu_D)}{C_3(M)}<1.03745
\]

for \(M/2<\mu_D<2M\).

Restoring the full Higgs doublet and colour multiplicities gives

\[
C_3^{\mathrm{full}}=1.58896\times10^{-4}\,\mathrm{GeV}^2,
\]

\[
|\Delta V_{Qa}^{(3)}|
=4.2902\times10^3\,\mathrm{GeV}^4,
\]

still only \(3.053\times10^{-12}\) of the thermal potential.

A separate warning survives: in the minimal one-loop high-scale run, the
Higgs quartic becomes negative. The matching calculation is stable, but
the displayed scalar sector is not yet a complete metastable UV
benchmark.

\section{Part XI - Explicit thermalisation and direct AMY
transport}\label{part-xi---explicit-thermalisation-and-direct-amy-transport}

\subsection{49. Reduced explicit collision
kernel}\label{reduced-explicit-collision-kernel}

The v1.3 calculation replaced a fitted BGK closure by a microscopic
discrete gain/loss operator for

\[
H,\quad D,\quad q,\quad g.
\]

It included six \(1\leftrightarrow2\) channel families, eight
\(2\leftrightarrow2\) families, Bose enhancement, Pauli blocking, Debye
screening, asymptotic thermal masses, an LPM formation-time reduction,
and exact discrete energy conservation.

The lattice contained 514 collinear transitions and 14,289 elastic
transitions. It reached \(99\%\) of the equilibrium entropy gain by
\(Tt=3920\), with detailed-balance residual \(8.57\times10^{-17}\) and
energy drift \(1.44\times10^{-15}\). The resulting hierarchy was

\[
\frac{\Gamma_{\mathrm{kin}}^{(99)}}{\Gamma_R}
=1.729\times10^6.
\]

Therefore homogeneous energy bookkeeping can adiabatically eliminate the
fast plasma without a phenomenological BGK rate
\citeproc{ref-AMY2002}{{[}45{]}}, \citeproc{ref-Boedeker2019}{{[}46{]}}.

\subsection{50. Direct AMY upgrade}\label{direct-amy-upgrade}

The v1.4 calculation directly solved the isotropic AMY transverse LPM
equation for the QCD splitting channels and replaced angle-averaged
elastic scattering by deterministic full-angle screened quadrature. At
the reheating benchmark

\[
\frac{M_D}{T}=0.01,
\qquad
\alpha_s=0.0393544,
\qquad
y_D=0.30,
\qquad
\frac pT=3,
\]

the dimensionless rates are

\[
\frac{\Gamma_{g\to gg}}T=0.261356,
\]

\[
\frac{\Gamma_{q\to gq}}T=0.204516,
\qquad
\frac{\Gamma_{D\to gD}}T=0.201417,
\]

\[
\frac{\Gamma_{g,\mathrm{elastic}}}T=0.106177,
\qquad
\frac{\Gamma_{q,\mathrm{elastic}}}T=0.0119151,
\]

and

\[
\boxed{
\frac{\Gamma_{H\leftrightarrow qD}}T
=3.18544\times10^{-4}.
}
\]

The QCD rates are fast. The slow mode is the scalar-Yukawa conversion
\(H\leftrightarrow qD\). At

\[
T_0=1.002\times10^8\,\mathrm{GeV},
\]

the kinetic rate is

\[
\Gamma_{\mathrm{kin}}=3.1918\times10^4\,\mathrm{GeV}.
\]

Compared with

\[
\Gamma_R=1.47850065\times10^{-2}\,\mathrm{GeV},
\]

one obtains

\[
\boxed{
\frac{\Gamma_{\mathrm{kin}}}{\Gamma_R}
=2.1588\times10^6.
}
\]

Even a conservative factor-of-two normalisation uncertainty leaves
\(1.0794\times10^6\). The correction to the desired reheating branch is
negligible.

\begin{figure}
\centering
\includegraphics[width=0.84\textwidth,height=\textheight]{figures/kinetic_rate_hierarchy.png}
\caption{The slow portal conversion remains more than six orders of
magnitude faster than reheaton decay.}\label{fig:rates}
\end{figure}

\subsection{51. What v1.4 did and did not
solve}\label{what-v1.4-did-and-did-not-solve}

The direct AMY calculation substantially closes the QCD kinetic problem.
It does not yet contain the complete electroweak/Yukawa LPM structure of
\(H\leftrightarrow qD\). The outstanding calculation requires:

\begin{enumerate}
\def\labelenumi{\arabic{enumi}.}
\tightlist
\item
  the complete helicity and Higgs-doublet source structure;
\item
  simultaneous \(SU(3)_c\), \(SU(2)_L\), and \(U(1)_Y\) soft collision
  kernels;
\item
  thermal widths and asymptotic masses;
\item
  exact matching to the scalar retarded self-energy;
\item
  insertion into a reduced gauge-covariant Schwinger-Keldysh model.
\end{enumerate}

Only after that does a full \(3+1\) dimensional non-Abelian
2PI/Kadanoff-Baym computation become the rational next expenditure. The
QCD sector is already far too fast to be the weak link.

\section{Part XII - Integrated interpretation, novelty, and
falsification}\label{part-xii---integrated-interpretation-novelty-and-falsification}

\subsection{52. What ``time emerges'' means
here}\label{what-time-emerges-means-here}

The final hierarchy is

\[
\boxed{
\begin{aligned}
\text{causal/conformal structure}
&\longrightarrow\text{before/after and null directions},\\
\text{nonparallel fields or phase gradients}
&\longrightarrow\text{local timelike energy directions},\\
\text{generated universal scale }\chi
&\longrightarrow\text{metric normalisation},\\
\text{material spectral gaps}
&\longrightarrow\text{stable relative phases},\\
\text{phase counting}
&\longrightarrow\text{operational duration}.
\end{aligned}}
\]

The programme derives a chronometric calibration conditional on a
conformal causal structure and physical quantum dynamics. It does not
derive:

\begin{itemize}
\tightlist
\item
  events or causal order;
\item
  the differentiable manifold;
\item
  Einstein's field equations;
\item
  a universal global present;
\item
  the thermodynamic arrow;
\item
  every quantum interaction;
\item
  a UV-complete theory of gravity.
\end{itemize}

The strongest defensible sentence is:

\begin{quote}
Causality may be primitive; proper duration can be the universally
calibrated phase relation generated by interacting fields.
\end{quote}

\subsection{53. Non-circularity}\label{non-circularity}

If \(\chi\) is a pure Weyl compensator, the construction can be
interpreted as gauge-invariant bookkeeping. No local dimensionless
experiment detects a universal rescaling of all masses. The result is
conceptually clean but may not constitute new dynamics.

If \(\chi\) is a physical condensate or dilaton, the theory becomes
stronger only if it supplies at least one of:

\begin{enumerate}
\def\labelenumi{\arabic{enumi}.}
\tightlist
\item
  a microscopic derivation of the scale;
\item
  gravitational or cosmological dynamics controlled by it;
\item
  controlled nonuniversal residuals;
\item
  a relation among clock, fifth-force, collider, or cosmological
  observables.
\end{enumerate}

The QCD defect, the \(Z_6\) protection, and the clock/EP consistency
relation are the project's attempt to meet this requirement.

\subsection{54. Novelty boundary}\label{novelty-boundary}

The following are established and should not be claimed as novel:

\begin{itemize}
\tightlist
\item
  null-pair and null-tetrad algebra;
\item
  reciprocal null scaling as Lorentz boost freedom;
\item
  timelike total momentum from nonparallel massless momenta;
\item
  causal order determining conformal structure under suitable
  conditions;
\item
  clocks selecting a conformal representative;
\item
  mimetic conformal maps;
\item
  two-superfluid entrainment;
\item
  dimensional transmutation;
\item
  Higgs-dilaton and induced-scale models;
\item
  QCD threshold decoupling and the \(2/27\) coefficient;
\item
  cyclic \(Z_N\) suppression of an ultralight scalar potential;
\item
  asymmetric reheating and finite-density potentials;
\item
  Schwinger-Keldysh, 2PI/KB, AMY, LPM, and screened collision methods.
\end{itemize}

The candidate contribution is narrower:

\[
\boxed{
\begin{gathered}
\text{universal clock factorisation and chronometric shear}
+\text{three-sector scale locking}\\
+\text{QCD-transmitted }2/27\text{ defect}
+\text{protected }Z_6\text{ long-range mode}\\
+\text{state-selected exact-symmetry reheating}
+\text{transient matching and microscopic transport}.
\end{gathered}}
\]

No exact indexed match for that full conjunction was located in the
staged searches. That is a candidate paper claim, not an absolute
priority claim. Publication requires backward and forward citation
chaining, non-arXiv databases, and specialist review in mimetic gravity,
relativistic superfluids, varying constants, axion quality,
nonequilibrium QFT, and thermal field theory.

\subsection{55. Falsification routes}\label{falsification-routes}

The programme is falsifiable at several levels.

\subsubsection{Soft-null analogue or hidden-sector
spectrum}\label{soft-null-analogue-or-hidden-sector-spectrum}

The minimal phase EFT predicts

\[
c_\parallel^2=2c_\perp^2-1,
\qquad
\alpha_\parallel=4\alpha_\perp.
\]

Deviation indicates additional operators or falsifies the minimal
mediator completion.

\subsubsection{Clock-network rank}\label{clock-network-rank}

One nonuniversal scalar predicts a rank-one matrix of differential
clock-ratio responses. Rank two or higher falsifies the single-mode
model, though not the general chronometric-shear formalism.

\subsubsection{Clock/EP consistency
line}\label{clockep-consistency-line}

Under one long-range unscreened QCD-dominated scalar,

\[
\frac{\beta_{AB}}{\eta_{CD}}
=-\frac{\Delta K_{AB}}{\Delta Q_{\widehat m}^{\prime\,CD}}.
\]

Inconsistent clock and free-fall signals would rule out the benchmark
assumptions.

\subsubsection{Cosmological sign and
amplitude}\label{cosmological-sign-and-amplitude}

The state-selected reheating construction predicts a selected sign of
\(d_g\), sparse adjacent-sector dark radiation, and negligible
chronometric isocurvature in the strong attractor. A full perturbation
calculation could falsify this architecture.

\subsubsection{Kinetic hierarchy}\label{kinetic-hierarchy}

The model requires the visible and adjacent plasma to equilibrate much
faster than reheaton decay. Direct AMY transport currently supports a
hierarchy above \(10^6\). A complete electroweak/Yukawa result would
have to lower the portal rate by more than six orders of magnitude to
overturn the cascade approximation.

\subsection{56. Integrated acceptance
matrix}\label{integrated-acceptance-matrix}

\begin{figure}
\centering
\includegraphics[width=0.82\textwidth,height=\textheight]{figures/integrated_status_counts.png}
\caption{The present project contains genuine passes, useful conditional
results, explicit failures, and several large open
problems.}\label{fig:status}
\end{figure}

The machine-readable integrated matrix is supplied as
\texttt{data/integrated\_acceptance\_matrix.csv}. The manuscript's
central open items are:

\begin{enumerate}
\def\labelenumi{\arabic{enumi}.}
\tightlist
\item
  complete electroweak/Yukawa LPM for \(H\leftrightarrow qD\);
\item
  two-loop or higher RG-improved vacuum and tunnelling analysis of the
  full scalar sector;
\item
  complete reheating cascade with the repaired fermionic parent,
  rescattering, and gauge interactions;
\item
  full cosmological perturbations and likelihoods;
\item
  UV completion of the collective Wilson line and all discrete
  anomalies;
\item
  a positive invariant spectrum for any proposed Turok-36 completion;
\item
  a non-circular microscopic origin of gravity, if the ambition extends
  beyond duration calibration;
\item
  a publication-grade novelty audit.
\end{enumerate}

\section{Conclusion}\label{conclusion}

The project began with a seductive but invalid phrase: ``the photon's
perspective.'' Relativity does not permit such an observer. What
survives is more interesting. A null direction supplies causal
propagation without an internal clock. Two nonparallel null relations
can supply a timelike direction, but not an absolute duration scale.
Interactions can generate a common spectral scale. A family of clocks
defines one proper-time metric exactly when every local spectrum shares
that factor. The failure of this condition is a Weyl-invariant,
experimentally meaningful chronometric shear.

The first exact-null dynamics failed. The theory became healthier by
softening nullity and separating kinematics from scale generation. The
Higgs alone proved insufficient, so a deeper scale field was introduced.
QCD thresholds made the first nonuniversal defect calculable. A \(Z_6\)
orbit protected an ultralight mode. Finite-size analysis showed that
ordinary bodies do not screen it. Cosmology selected a branch only after
the parameter space moved away from the original Planck-scale benchmark.
State-selected reheating preserved vacuum symmetry, while
closed-time-path analysis showed why transient asymmetry need not become
a permanent spurion. A bosonic preheating portal failed and was replaced
by a fermionic cascade. Exact factorised matching, RG completion,
explicit collision kernels, and direct AMY transport progressively
reduced the remaining uncertainty.

The result is not a completed theory of time, gravity, or quantum
cosmology. It is a coherent and increasingly constrained theory of
\textbf{duration calibration}:

\[
\boxed{
\begin{gathered}
\text{null incidence}
+\text{interaction-generated universal scale}\\
+\text{spectral factorisation}
+\text{phase comparison}\\
\Longrightarrow
\text{operational proper duration}.
\end{gathered}
}
\]

The next decisive calculation is sharply localised. It is the complete
electroweak/Yukawa LPM problem for \(H\leftrightarrow qD\). That
calculation will determine the last meaningful kinetic normalisation
before the project earns the computational bloodletting of a full
non-Abelian \(3+1\) dimensional two-time treatment.

\newpage

\section{Appendix A - Core notation}\label{appendix-a---core-notation}

\begin{longtable}[]{@{}
  >{\raggedright\arraybackslash}p{(\columnwidth - 2\tabcolsep) * \real{0.5000}}
  >{\raggedright\arraybackslash}p{(\columnwidth - 2\tabcolsep) * \real{0.5000}}@{}}
\toprule\noalign{}
\begin{minipage}[b]{\linewidth}\raggedright
Symbol
\end{minipage} & \begin{minipage}[b]{\linewidth}\raggedright
Meaning
\end{minipage} \\
\midrule\noalign{}
\endhead
\bottomrule\noalign{}
\endlastfoot
\([g]\) & Lorentzian conformal class \\
\(k_a,\ell_a\) & crossed null phase gradients or null covectors \\
\(C=-\tfrac12k\cdot\ell\) & crossed-null invariant scale \\
\(T,R\) & mean and difference phases \\
\(h_{ab}=Cg_{ab}\) & singular crossed-phase/mimetic representative in
exact-null formulation \\
\(\chi\) & common Weyl-weight \(-1\) spectral scale \\
\(\widehat g_{\mu\nu}\) & operational chronometric metric \\
\(\omega_A\) & local spectrum of clock transition \(A\) \\
\(\mathcal S_{AB}\) & chronometric shear \(d\ln(\omega_A/\omega_B)\) \\
\(\delta_n\) & QCD lock defect \(d\ln(\Lambda_n/\chi)\) \\
\(\epsilon_j\) & threshold defect \(d\ln(\widehat m_j/\chi)\) \\
\(a=f_a\theta\) & canonical ratio mode \\
\(d_g\) & dimensionless scalar coupling to the low-energy QCD scale \\
\(x=a/f_a\) & compact phase variable \\
\(Q_k\) & transient selector orbit \\
\(X_k,Y_k,R_k\) & reheaton orbit fields/eigenstates \\
\(N_k\) & selector-gated fermionic parent orbit \\
\(H,D,q,g\) & Higgs, vectorlike fermion, light quark, gluon kinetic
species \\
\end{longtable}

\section{Appendix B - Consolidated benchmark
values}\label{appendix-b---consolidated-benchmark-values}

The machine-readable source is \texttt{data/benchmark\_parameters.json}.
Key values are reproduced here.

\begin{longtable}[]{@{}lr@{}}
\toprule\noalign{}
Benchmark & Value \\
\midrule\noalign{}
\endhead
\bottomrule\noalign{}
\endlastfoot
Reduced Planck mass & \(2.435\times10^{18}\) GeV \\
QCD threshold transmission & \(2/27\) \\
\(Z_6\) Planck benchmark range & \(347.4\) AU \\
Environmental conversion threshold \(q_{\mathrm{conv}}\) &
\(0.633135\) \\
Strong-attractor \(f_a\) & \(2.435\times10^{10}\) GeV \\
Strong-attractor \(\epsilon\) & \(2.70\times10^{-13}\) \\
Strong-attractor \(M\) & \(1.002\times10^6\) GeV \\
Strong-attractor \(d_g\) & \(10^{-6}\) \\
Adjacent-sector \(\Delta N_{\mathrm{eff}}\) & \(0.0289\) \\
Updated \(g_*\) & \(117.25\) \\
Updated reheaton width & \(1.47850065\times10^{-2}\) GeV \\
Updated adjacent branch \(B_5\) & \(0.00529888708\) \\
RG-improved \(\mathcal I_3\) & \(6.57973508149\) \\
Transient/thermal potential ratio & \(3.053\times10^{-12}\) \\
Direct-AMY kinetic rate & \(3.1918\times10^4\) GeV \\
Direct-AMY hierarchy & \(2.1588\times10^6\) \\
\end{longtable}

\section{Appendix C - Version-by-version correction
ledger}\label{appendix-c---version-by-version-correction-ledger}

\begin{longtable}[]{@{}
  >{\raggedright\arraybackslash}p{(\columnwidth - 4\tabcolsep) * \real{0.3333}}
  >{\raggedright\arraybackslash}p{(\columnwidth - 4\tabcolsep) * \real{0.3333}}
  >{\raggedright\arraybackslash}p{(\columnwidth - 4\tabcolsep) * \real{0.3333}}@{}}
\toprule\noalign{}
\begin{minipage}[b]{\linewidth}\raggedright
Version
\end{minipage} & \begin{minipage}[b]{\linewidth}\raggedright
Retained result
\end{minipage} & \begin{minipage}[b]{\linewidth}\raggedright
Correction introduced
\end{minipage} \\
\midrule\noalign{}
\endhead
\bottomrule\noalign{}
\endlastfoot
v0.1 & Exact crossed-phase algebra and product/ratio split & Homogeneous
stability was not field-theory stability. \\
v0.2 & Soft-null EFT and mediator completion & Exact-null map is
mimetic; constrained action strongly coupled transversely. \\
v0.3 & Universal factorisation theorem and chronometric shear & Ordinary
Higgs is insufficient as sole common scale. \\
v0.4 & Conventional hidden scale generator & Turok-36 moved from
assumption to optional candidate under hard tests. \\
v0.5 & All-orders threshold recursion and \(2/27\) & Minimal long-range
defect is radiatively unnatural. \\
v0.6 & \(Z_6\) protection & Global scalar quality and finite-density
effects become decisive. \\
v0.7 & No environmental screening & Cosmological homogeneous selection
remains separate. \\
v0.8 & Low-\(f_a\) strong attractor & Original Planck benchmark
overcloses and does not select a vacuum. \\
v0.9 & State-selected exact-symmetry reheating & Naive branching and
preheating still unresolved. \\
v1.0 & Vacuum/state CTP separation through two loops & Exact \(Z_6\)
alone is insufficient without shift spurion counting. \\
v1.1 & Three-loop transient topology and fermionic repair & Direct
bosonic portal fails generically. \\
v1.2 & Exact factorised matching and nonlinear cascade &
\(\tan\theta=1/16\) fails; corrected branch required. \\
v1.3 & RG completion and explicit collision kernel & Fixed-order scale
excursion and BGK closure removed. \\
v1.4 & Direct QCD AMY transport & Electroweak/Yukawa LPM, not QCD, is
the remaining kinetic bottleneck. \\
\end{longtable}

\section{Appendix D - Reproducibility and
provenance}\label{appendix-d---reproducibility-and-provenance}

This consolidation is based on:

\begin{enumerate}
\def\labelenumi{\arabic{enumi}.}
\tightlist
\item
  the current-conversation chat snapshot stored in the \texttt{sources}
  directory;
\item
  the staged manuscript, source, code, result, matrix, and figure
  records from v0.1-v1.4 located in the project File Library;
\item
  the equations, benchmarks, caveats, and correction history explicitly
  reported in those staged notes;
\item
  a source ledger that preserves the cited literature, corrected
  metadata, and source-role audit notes.
\end{enumerate}

Historical sandbox links from earlier chats are not assumed to be
durable local paths. The archive therefore contains a complete
\textbf{consolidated} manuscript, bibliography, source ledger,
verification suite, benchmark data, figures, and an inventory of the
historical filenames. It does not falsely claim to contain
byte-identical copies of every ephemeral historical ZIP or figure. See
\texttt{data/historical\_artifact\_inventory.csv} and
\texttt{SOURCE\_PROVENANCE.md}.

The verification suite checks the main algebraic and numerical
identities reproduced in this manuscript. It is not a substitute for
peer review, independent code, lattice gauge theory, or the unresolved
high-performance calculations.

\section*{References}\label{bibliography}
\addcontentsline{toc}{section}{References}

\phantomsection\label{refs}
\begin{CSLReferences}{0}{0}
\bibitem[\citeproctext]{ref-EPS1972}
\CSLLeftMargin{{[}1{]} }%
\CSLRightInline{J. Ehlers, F. A. E. Pirani, and A. Schild, {`The
geometry of free fall and light propagation'}. GRG reprint 2012, 1972.}

\bibitem[\citeproctext]{ref-HKM1976}
\CSLLeftMargin{{[}2{]} }%
\CSLRightInline{S. W. Hawking, A. R. King, and P. J. McCarthy, {`A new
topology for curved space-time which incorporates the causal,
differential, and conformal structures'}. J. Math. Phys. 17, 174, 1976.}

\bibitem[\citeproctext]{ref-Malament1977}
\CSLLeftMargin{{[}3{]} }%
\CSLRightInline{D. B. Malament, {`The class of continuous timelike
curves determines the topology of spacetime'}. J. Math. Phys. 18, 1399,
1977.}

\bibitem[\citeproctext]{ref-Braun2025}
\CSLLeftMargin{{[}4{]} }%
\CSLRightInline{M. Braun, {`Spacetime reconstruction by order and
number'}. 2025. Available: \url{https://arxiv.org/abs/2507.01907}}

\bibitem[\citeproctext]{ref-Chetyrkin1997}
\CSLLeftMargin{{[}5{]} }%
\CSLRightInline{K. G. Chetyrkin, B. A. Kniehl, and M. Steinhauser,
{`Decoupling relations to o(alpha\_s\^{}3) and their connection to
low-energy theorems'}. 1997. Available:
\url{https://arxiv.org/abs/hep-ph/9708255}}

\bibitem[\citeproctext]{ref-DamourDonoghue2010}
\CSLLeftMargin{{[}6{]} }%
\CSLRightInline{T. Damour and J. F. Donoghue, {`Equivalence principle
violations and couplings of a light dilaton'}. 2010. Available:
\url{https://arxiv.org/abs/1007.2792}}

\bibitem[\citeproctext]{ref-Gourgoulhon2006}
\CSLLeftMargin{{[}7{]} }%
\CSLRightInline{E. Gourgoulhon and J. L. Jaramillo, {`A 3+1 perspective
on null hypersurfaces and isolated horizons'}. 2006. Available:
\url{https://arxiv.org/abs/gr-qc/0503113}}

\bibitem[\citeproctext]{ref-Mimetic2013}
\CSLLeftMargin{{[}8{]} }%
\CSLRightInline{A. H. Chamseddine and V. Mukhanov, {`Mimetic dark
matter'}. 2013. Available: \url{https://arxiv.org/abs/1308.5410}}

\bibitem[\citeproctext]{ref-Superfluid2015}
\CSLLeftMargin{{[}9{]} }%
\CSLRightInline{A. Haber, A. Schmitt, and S. Stetina, {`Instabilities in
relativistic two-component (super)fluids'}. 2015. Available:
\url{https://arxiv.org/abs/1510.01982}}

\bibitem[\citeproctext]{ref-Adams2006}
\CSLLeftMargin{{[}10{]} }%
\CSLRightInline{A. Adams, N. Arkani-Hamed, S. Dubovsky, A. Nicolis, and
R. Rattazzi, {`Causality, analyticity and an IR obstruction to UV
completion'}. 2006. Available:
\url{https://arxiv.org/abs/hep-th/0602178}}

\bibitem[\citeproctext]{ref-Atkins2012}
\CSLLeftMargin{{[}11{]} }%
\CSLRightInline{M. Atkins and X. Calmet, {`Bounds on the nonminimal
coupling of the higgs boson to gravity'}. 2012. Available:
\url{https://arxiv.org/abs/1211.0281}}

\bibitem[\citeproctext]{ref-Shaposhnikov2009}
\CSLLeftMargin{{[}12{]} }%
\CSLRightInline{M. Shaposhnikov and D. Zenhausern, {`Scale invariance,
unimodular gravity and dark energy'}. 2009. Available:
\url{https://arxiv.org/abs/0809.3395}}

\bibitem[\citeproctext]{ref-Foot2007}
\CSLLeftMargin{{[}13{]} }%
\CSLRightInline{R. Foot, A. Kobakhidze, K. L. McDonald, and R. R.
Volkas, {`A solution to the hierarchy problem from an almost decoupled
hidden sector within a classically scale invariant theory'}. 2007.
Available: \url{https://arxiv.org/abs/0709.2750}}

\bibitem[\citeproctext]{ref-HiddenSU22013}
\CSLLeftMargin{{[}14{]} }%
\CSLRightInline{C. D. Carone and R. Ramos, {`Classical scale-invariance,
the electroweak scale and vector dark matter'}. 2013. Available:
\url{https://arxiv.org/abs/1307.8428}}

\bibitem[\citeproctext]{ref-CurvedSM2018}
\CSLLeftMargin{{[}15{]} }%
\CSLRightInline{T. Markkanen, S. Nurmi, A. Rajantie, and S. Stopyra,
{`The 1-loop effective potential for the standard model in curved
spacetime'}. 2018. Available: \url{https://arxiv.org/abs/1804.02020}}

\bibitem[\citeproctext]{ref-ScaleFifthForce2016}
\CSLLeftMargin{{[}16{]} }%
\CSLRightInline{P. G. Ferreira, C. T. Hill, and G. G. Ross, {`No fifth
force in a scale invariant universe'}. 2016. Available:
\url{https://arxiv.org/abs/1612.03157}}

\bibitem[\citeproctext]{ref-ScaleFifthForce2021}
\CSLLeftMargin{{[}17{]} }%
\CSLRightInline{E. J. Copeland, P. Millington, and S. S. Muñoz, {`Fifth
forces and broken scale symmetries in the jordan frame'}. 2021.
Available: \url{https://arxiv.org/abs/2111.06357}}

\bibitem[\citeproctext]{ref-BoyleTurok2022}
\CSLLeftMargin{{[}18{]} }%
\CSLRightInline{L. Boyle and N. Turok, {`Cancelling the vacuum energy
and weyl anomaly in the standard model with dimension-zero scalar
fields'}. 2022. Available: \url{https://arxiv.org/abs/2110.06258}}

\bibitem[\citeproctext]{ref-BatemanTurok2026}
\CSLLeftMargin{{[}19{]} }%
\CSLRightInline{S. Bateman and N. Turok, {`Escape from ostrogradsky via
hidden ghost parity'}. 2026. Available:
\url{https://arxiv.org/abs/2607.00096}}

\bibitem[\citeproctext]{ref-ClineHell2026}
\CSLLeftMargin{{[}20{]} }%
\CSLRightInline{J. M. Cline and A. Hell, {`Pathologies of dimension-zero
scalar fields'}. 2026. Available:
\url{https://arxiv.org/abs/2603.05683}}

\bibitem[\citeproctext]{ref-ClockData2023}
\CSLLeftMargin{{[}21{]} }%
\CSLRightInline{N. Sherrill \emph{et al.}, {`Analysis of atomic-clock
data to constrain variations of fundamental constants'}. 2023.
Available: \url{https://arxiv.org/abs/2302.04565}}

\bibitem[\citeproctext]{ref-ClockReview2025}
\CSLLeftMargin{{[}22{]} }%
\CSLRightInline{A. Kawasaki, {`Quantum sensing using atomic clocks for
nuclear and particle physics'}. 2025. Available:
\url{https://arxiv.org/abs/2411.19424}}

\bibitem[\citeproctext]{ref-MICROSCOPE2022}
\CSLLeftMargin{{[}23{]} }%
\CSLRightInline{P. Touboul \emph{et al.}, {`MICROSCOPE mission: Final
results of the test of the equivalence principle'}. 2022. Available:
\url{https://arxiv.org/abs/2209.15487}}

\bibitem[\citeproctext]{ref-DasHook2020}
\CSLLeftMargin{{[}24{]} }%
\CSLRightInline{S. Das and A. Hook, {`Non-linearly realized discrete
symmetries'}. 2020. Available: \url{https://arxiv.org/abs/2006.10767}}

\bibitem[\citeproctext]{ref-Brzeminski2020}
\CSLLeftMargin{{[}25{]} }%
\CSLRightInline{D. Brzeminski, Z. Chacko, A. Dev, and A. Hook, {`A
time-varying fine structure constant from naturally ultralight dark
matter'}. 2020. Available: \url{https://arxiv.org/abs/2012.02787}}

\bibitem[\citeproctext]{ref-DiLuzio2021}
\CSLLeftMargin{{[}26{]} }%
\CSLRightInline{L. D. Luzio, B. Gavela, P. Quilez, and A. Ringwald, {`An
even lighter QCD axion'}. 2021. Available:
\url{https://arxiv.org/abs/2102.00012}}

\bibitem[\citeproctext]{ref-Hor2026}
\CSLLeftMargin{{[}27{]} }%
\CSLRightInline{S. Hor, Y. Nakai, M. Suzuki, and J. Xu, {`Deconstructing
the extra-dimensional axion'}. 2026. Available:
\url{https://arxiv.org/abs/2606.02728}}

\bibitem[\citeproctext]{ref-Chameleon2008}
\CSLLeftMargin{{[}28{]} }%
\CSLRightInline{T. Tamaki and S. Tsujikawa, {`Revisiting chameleon
gravity---thin-shells and no-shells with appropriate boundary
conditions'}. 2008. Available: \url{https://arxiv.org/abs/0808.2284}}

\bibitem[\citeproctext]{ref-FiniteDensity2024}
\CSLLeftMargin{{[}29{]} }%
\CSLRightInline{O. F. Ramadan, J. Sakstein, and D. Croon, {`Cosmology
and astrophysics of CP-violating axions'}. 2024. Available:
\url{https://arxiv.org/abs/2408.02294}}

\bibitem[\citeproctext]{ref-GoldsteinHill2026}
\CSLLeftMargin{{[}30{]} }%
\CSLRightInline{S. Goldstein and J. C. Hill, {`A 2\% determination of
n\_eff from primordial element abundance, CMB and BAO measurements'}.
2026. Available: \url{https://arxiv.org/abs/2603.13226}}

\bibitem[\citeproctext]{ref-Borsanyi2016}
\CSLLeftMargin{{[}31{]} }%
\CSLRightInline{Sz. Borsanyi \emph{et al.}, {`Lattice QCD for
cosmology'}. 2016. Available: \url{https://arxiv.org/abs/1606.07494}}

\bibitem[\citeproctext]{ref-Arakawa2026}
\CSLLeftMargin{{[}32{]} }%
\CSLRightInline{J. Arakawa \emph{et al.}, {`Probing ultralight dark
matter at the mega-planck scale with the thorium nuclear clock'}. 2026.
Available: \url{https://arxiv.org/abs/2602.16804}}

\bibitem[\citeproctext]{ref-Thorium2026}
\CSLLeftMargin{{[}33{]} }%
\CSLRightInline{L. T. D. Col \emph{et al.}, {`A thorium-229 optical
nuclear clock with feedback loop'}. 2026. Available:
\url{https://arxiv.org/abs/2606.04997}}

\bibitem[\citeproctext]{ref-Delaunay2026}
\CSLLeftMargin{{[}34{]} }%
\CSLRightInline{C. Delaunay, S. J. Lee, Y. Yin, and B. Yu, {`Natural
phantom crossing from axion-WIMP interactions'}. 2026. Available:
\url{https://arxiv.org/abs/2607.28721}}

\bibitem[\citeproctext]{ref-ThermalPressure2021}
\CSLLeftMargin{{[}35{]} }%
\CSLRightInline{J.-L. Kneur, M. B. Pinto, and T. E. Restrepo,
{`Renormalization group improved pressure for hot and dense quark
matter'}. 2021. Available: \url{https://arxiv.org/abs/2101.08240}}

\bibitem[\citeproctext]{ref-CollinsHolman2005}
\CSLLeftMargin{{[}36{]} }%
\CSLRightInline{H. Collins and R. Holman, {`Renormalization of initial
conditions and the trans-planckian problem of inflation'}. 2005.
Available: \url{https://arxiv.org/abs/hep-th/0501158}}

\bibitem[\citeproctext]{ref-Berges2005}
\CSLLeftMargin{{[}37{]} }%
\CSLRightInline{J. Berges, S. Borsanyi, U. Reinosa, and J. Serreau,
{`Nonperturbative renormalization for 2PI effective action techniques'}.
2005. Available: \url{https://arxiv.org/abs/hep-ph/0503240}}

\bibitem[\citeproctext]{ref-Tranberg2024}
\CSLLeftMargin{{[}38{]} }%
\CSLRightInline{A. Tranberg and G. Ungersbaeck, {`Four results on
out-of-equilibrium 2PI simulations in 3+1 dimensions'}. 2024. Available:
\url{https://arxiv.org/abs/2409.06398}}

\bibitem[\citeproctext]{ref-Lindner2007}
\CSLLeftMargin{{[}39{]} }%
\CSLRightInline{M. Lindner and M. M. Muller, {`Comparison of boltzmann
kinetics with quantum dynamics for a chiral yukawa model far from
equilibrium'}. 2007. Available: \url{https://arxiv.org/abs/0710.2917}}

\bibitem[\citeproctext]{ref-Bhattacharya2022}
\CSLLeftMargin{{[}40{]} }%
\CSLRightInline{S. Bhattacharya, N. Joshi, and S. Kaushal, {`Decoherence
and entropy generation in an open quantum scalar-fermion system with
yukawa interaction'}. 2022. Available:
\url{https://arxiv.org/abs/2206.15045}}

\bibitem[\citeproctext]{ref-Dufaux2006}
\CSLLeftMargin{{[}41{]} }%
\CSLRightInline{J. F. Dufaux, G. N. Felder, L. Kofman, M. Peloso, and D.
Podolsky, {`Preheating with trilinear interactions: Tachyonic
resonance'}. 2006. Available:
\url{https://arxiv.org/abs/hep-ph/0602144}}

\bibitem[\citeproctext]{ref-GreeneKofman1998}
\CSLLeftMargin{{[}42{]} }%
\CSLRightInline{P. B. Greene and L. Kofman, {`Preheating of fermions'}.
1998. Available: \url{https://arxiv.org/abs/hep-ph/9807339}}

\bibitem[\citeproctext]{ref-Byakti2017}
\CSLLeftMargin{{[}43{]} }%
\CSLRightInline{P. Byakti, D. Ghosh, and T. Sharma, {`Note on gauge and
gravitational anomalies of discrete z\_n symmetries'}. 2017. Available:
\url{https://arxiv.org/abs/1707.03837}}

\bibitem[\citeproctext]{ref-Martin2001}
\CSLLeftMargin{{[}44{]} }%
\CSLRightInline{S. P. Martin, {`Two-loop effective potential for a
general renormalizable theory and softly broken supersymmetry'}. 2001.
Available: \url{https://arxiv.org/abs/hep-ph/0111209}}

\bibitem[\citeproctext]{ref-AMY2002}
\CSLLeftMargin{{[}45{]} }%
\CSLRightInline{P. Arnold, G. D. Moore, and L. G. Yaffe, {`Effective
kinetic theory for high temperature gauge theories'}. 2002. Available:
\url{https://arxiv.org/abs/hep-ph/0209353}}

\bibitem[\citeproctext]{ref-Boedeker2019}
\CSLLeftMargin{{[}46{]} }%
\CSLRightInline{D. Boedeker and D. Schroeder, {`Equilibration of
right-handed electrons'}. 2019. Available:
\url{https://arxiv.org/abs/1902.07220}}

\end{CSLReferences}

\end{document}
